\subsection{The case of groups without arbitrarily small finite subgroups}

THEOREM 2. --- Let G be a locally compact group satisfying the following condition:

\begin{enumerate}
\def\labelenumi{(\Alph{enumi})}
\setcounter{enumi}{11}
\tightlist
\item
  There exists a neighborhood of \(e\) in \(G\) that contains no finite subgroup of \(G\) not reduced to \(e\) .
\end{enumerate}

The following properties then hold:

\begin{enumerate}
\def\labelenumi{(\roman{enumi})}
\item
  The set D of discrete subgroups of G is locally closed in \(\Sigma\) (which is equivalent to saying that it is locally compact).
\item
  For a closed subset \(\mathbf{A}\) of \(\mathbf{D}\) to be compact, it is necessary and sufficient that there exist a neighborhood \(\mathbf{U}\) of \(e\) in \(\mathbf{G}\) such that \(\mathbf{H} \cap \mathbf{U} = \{e\}\) for every subgroup \(\mathbf{H} \in \mathbf{A}\) .
\item
  If in addition G is generated by a compact neighborhood of e, then the set \(D_{c}\) of discrete subgroups H of G such that G/H is compact is locally closed in \(\Sigma\) (hence is locally compact).
\end{enumerate}

We have \(\mathbf{D}_c = \mathbf{D} \cap \Sigma_c^0\) , therefore (iii) is a consequence of (i), and Th. 1 (ii) of No.~3.

With the notations of No.~3, Prop. 7, it suffices, for proving (i) and (ii), to prove that:

PROPOSITION 8. --- The canonical bijection of N onto D is a homeomorphism.

Now, if \(\Gamma_d\) is the set of Haar measures on the discrete subgroups of G, then D is canonically homeomorphic to the space of orbits of the group of homotheties in \(\Gamma_d\) with ratio \(>0\) (GT, I, §5, No.~2, Prop. 4). It therefore suffices to prove that the canonical mapping \(\alpha \mapsto (\alpha(\{e\}), \alpha/\alpha(\{e\}))\) of \(\Gamma_d\) onto \(\mathbf{R}_+^* \times \mathbf{N}\) is a homeomorphism, which will result from the following lemma:

Lemma 4. --- If the group G satisfies the condition (L), the mapping \(\alpha \mapsto \alpha (\{e\})\) of \(\Gamma_d\) into \(\mathbf{R}_+^*\) is vaguely continuous.

Let us consider a measure \(\alpha \in \Gamma_d\) ; let \(V_0\) be a relatively compact open neighborhood of \(e\) in \(G\) such that \(H_\alpha \cap V_0 = \{e\}\) and such that there exists no finite subgroup of \(G\) contained in \(V_0\) and not reduced to \(e\) . Let \(V\) be a symmetric compact neighborhood of \(e\) such that \(V^3 \subset V_0\) , and let \(U\) be a symmetric neighborhood of \(e\) such that \(U^2 \subset V\) . Let \(\varphi\) (resp. \(\psi\) ) be a function in \(\mathcal{K}_+(G)\) , with values in \([0,1]\) , equal to 1 on \(V^3\) (resp. at the point \(e\) ) and with support contained in \(V_0\) (resp. in \(U\) ). The set of measures \(\beta \in \Gamma_d\) such that \(|\beta(\varphi) - \alpha(\varphi)| \leqslant \varepsilon\) and \(|\beta(\psi) - \alpha(\psi)| \leqslant \varepsilon\) is a neighborhood \(W\) of \(\alpha\) . We propose to show that, provided \(\varepsilon\) is taken to be sufficiently small, \(H_\beta \cap V = \{e\}\) for every \(\beta \in W\) ; it will then follow that \(\beta(\psi) = \beta(\{e\})\) , hence that \(|\beta(\{e\}) - \alpha(\{e\})| \leqslant \varepsilon\) , which will prove the lemma.

It will suffice to show that, for \(\beta \in \mathbf{W}\) ,

\[
(\mathrm{V} ^ {2} - \mathrm{V}) \cap \mathrm{H} _ {\beta} = \varnothing .\tag{8}
\]

For, suppose that this point is established: then, for \(x\) and \(y\) in \(\mathbf{V} \cap \mathbf{H}_{\beta}\) , one has \(xy^{-1} \in \mathbf{V}^2 \cap \mathbf{H}_{\beta}\) ; but, by virtue of (8), this implies \(xy^{-1} \in \mathbf{V} \cap \mathbf{H}_{\beta}\) ; in other words, \(\mathbf{V} \cap \mathbf{H}_{\beta}\) is a subgroup of \(\mathbf{G}\) , which is obviously discrete and compact, hence finite; but then, by the choice of \(\mathbf{V}_0\) , this implies that indeed \(\mathbf{V} \cap \mathbf{H}_{\beta} = \{e\}\) .

Let us argue by contradiction and so assume that there exists a point \(z\) of \(\mathbf{V}^2 -\mathbf{V}\) that belongs to \(\mathbf{H}_{\beta}\) ; by the choice of U and V, we have \(\psi (sz^{-1}) + \psi (s)\leqslant \varphi (s)\) in G, the relation \(z\notin \mathrm{U}^2\) implying \(\mathrm{U}z\cap \mathrm{U} = \emptyset\) . Since

\[
\int \psi (s z ^ {- 1}) d \beta (s) = \int \psi (s) d \beta (s),
\]

it follows that \(2\beta (\psi)\leqslant \beta (\varphi)\leqslant \alpha (\varphi) + \varepsilon\) ; but we also have

\[
\beta (\psi) \geqslant \alpha (\psi) - \varepsilon ,
\]

and by construction \(\alpha(\varphi) = \alpha(\psi) = \alpha(\{e\})\) . We thus arrive at a contradiction by taking \(\varepsilon < \alpha(\{e\})/3\) . Q.E.D.

In graphic terms, a group G satisfying the condition (L) is said to not have arbitrarily small finite subgroups. *It can be shown that every Lie group satisfies the condition (L); but this condition is not characteristic of Lie groups; for example, the multiplicative group of \(p\) -adic integers congruent to 1 mod \(p\) satisfies (L).*

\subsection{The case of abelian groups}

Let G be a locally compact group, \(N \subset \Gamma^{0}\) the subspace of normalized Haar measures on the discrete subgroups of G, and \(N_{c}\) the subset of N corresponding to the discrete subgroups H of G such that G/H is compact; thus \(N_{c} = N \cap \Gamma_{c}^{0}\) in the notations of No.~3, Prop. 6; and if the group G is generated by a compact neighborhood of e, it follows from No.~3, Prop. 6 that \(N_{c}\) is open in N (hence is locally compact by No.~3, Cor. of Prop. 7) and that the restriction to \(N_{c}\) of the mapping \(\alpha \mapsto \|\mu_{\alpha}\|\) is vaguely continuous.

PROPOSITION 9. --- Let G be a locally compact abelian group, generated by a compact neighborhood of e. For a subset A of \(N_{c}\) to be relatively compact in \(N_{c}\) , it is necessary and sufficient that it satisfy the following two conditions:

\begin{enumerate}
\def\labelenumi{(\roman{enumi})}
\item
  There exists an open neighborhood \(\mathbf{U}\) of \(e\) in \(\mathbf{G}\) such that \(\mathrm{H}_{\alpha} \cap \mathrm{U} = \{e\}\) for all \(\alpha \in \mathbf{A}\) .
\item
  There exists a constant \(k\) such that \(\mu_{\alpha}(\mathrm{G} / \mathrm{H}_{\alpha}) \leqslant k\) for all \(\alpha \in \mathbf{A}\) . If \(\mathbf{A} \subset \mathbf{N}_c\) is relatively compact in \(\mathbf{N}_c\) , it is a fortiori so in \(\mathbf{N}\) , and the necessity of the conditions (i) and (ii) therefore follows from No.~3, Props. 6 and 7 (without assuming \(\mathbf{G}\) to be abelian). Conversely, suppose that \(\mathbf{A} \subset \mathbf{N}_c\) satisfies these conditions; if \(\overline{\mathbf{A}}\) is the closure of \(\mathbf{A}\) in \(\mathbf{N}\) , then \(\overline{\mathbf{A}}\) is compact by virtue of No.~3, Prop. 7; moreover, since \(\alpha \mapsto \| \mu_{\alpha} \|\) is lower semi-continuous on \(\Gamma^0\) for the vague topology (No.~2, Prop. 4), the condition (ii) implies that one also has \(\| \mu_{\alpha} \| \leqslant k\) for all \(\alpha \in \overline{\mathbf{A}}\) . Now, since \(\mathbf{G}\) is abelian, \(\mu_{\alpha} = \mu / \alpha\) is a Haar measure on the group \(G / H_{\alpha}\) , and \(G / H_{\alpha}\) is therefore compact for every \(\alpha \in \overline{\mathbf{A}}\) (Ch. VII, §1, No.~2, Prop. 2). This means that \(\overline{\mathbf{A}} \subset N_c\) , thus \(\mathbf{A}\) is relatively compact in \(N_c\) .
\end{enumerate}

COROLLARY. --- Let G be a locally compact abelian group, generated by a compact neighborhood of e and satisfying the condition (L) of No.~4.

Let \(D_c\) be the set of discrete subgroups H of G such that G/H is compact, and, for every \(H \in D_c\) , let \(v(H)\) be the total mass \(\mu_\alpha(G/H)\) , where \(\mu_\alpha\) is the quotient measure of \(\mu\) by the normalized Haar measure \(\alpha\) of H. For a subset A of the space \(D_c\) to be relatively compact in \(D_c\) , it is necessary and sufficient that it satisfy the following two conditions:

\begin{enumerate}
\def\labelenumi{(\roman{enumi})}
\item
  There exists an open neighborhood U of e in G such that \(\mathbf{H} \cap \mathbf{U} = \{e\}\) for all \(\mathbf{H} \in \mathbf{A}\) .
\item
  There exists a constant \(k\) such that \(v(\mathbf{H}) \leqslant k\) for all \(\mathbf{H} \in \mathbf{A}\) .
\end{enumerate}

Taking into account Prop. 9, this follows at once from the fact that \(D_{c}\) is the image of \(N_{c}\) under the canonical bijection of N onto D, and the fact that, under the hypotheses made, this bijection is a homeomorphism (No.~4, Prop. 8).

Example. --- Let us take \(G = R^{n}\) and for \(\mu\) the Lebesgue measure; all of the hypotheses of the Cor. of Prop. 9 are satisfied. The discrete subgroups H of G such that G/H is compact are none other than the discrete subgroups of rank n (GT, VII, §1, No.~1, Th. 1); such a subgroup H is generated by a basis \((a_{i})_{1 \leqslant i \leqslant n}\) of \(R^{n}\) , and

\[
v (\mathbf {H}) = | \det (a _ {1}, \dots , a _ {n}) |
\]

(the determinant being taken with respect to the canonical basis of \(\mathbf{R}^n\) ) (Ch. VII, §2, No.~10, Th. 4). The space \(\mathbf{D}_c\) can be interpreted here in the following way: every subgroup \(\mathbf{H} \in \mathbf{D}_c\) is the transform \(g \cdot \mathbf{Z}^n\) of the subgroup \(\mathbf{Z}^n\) by an element \(g \in \mathbf{GL}(n, \mathbf{R})\) , and the subgroup of \(\mathbf{GL}(n, \mathbf{R})\) leaving \(\mathbf{Z}^n\) stable may be identified with \(\mathbf{GL}(n, \mathbf{Z})\) . Consequently \(\mathbf{D}_c\) may be canonically identified, as a (non-topological) homogeneous space, with \(\mathbf{GL}(n, \mathbf{R}) / \mathbf{GL}(n, \mathbf{Z})\) . On the other hand, \(\mathbf{GL}(n, \mathbf{R})\) operates continuously in \(\mathbf{R}^n\) , hence also in \(\mathcal{M}_{+}(\mathbf{R}^{n})\) for the vague topology (§3, No.~3, Prop. 13), hence in the subspace \(\mathrm{N}_c\) of \(\mathcal{M}_{+}(\mathbf{R}^{n})\) ; moreover, the canonical homeomorphism (No.~4, Prop. 8) of \(\mathrm{N}_c\) onto \(\mathrm{D}_c\) is compatible with the laws of operation of \(\mathbf{GL}(n, \mathbf{R})\) . Since \(\mathbf{GL}(n, \mathbf{R})\) is countable at infinity and \(\mathrm{D}_c\) is locally compact, the bijection of \(\mathbf{GL}(n, \mathbf{R}) / \mathbf{GL}(n, \mathbf{Z})\) onto \(\mathrm{D}_c\) defined above is a homeomorphism (Ch. VII, App. I, Lemma 2). The Cor. of Prop. 9 therefore gives a criterion for compactness in the homogeneous space \(\mathbf{GL}(n, \mathbf{R}) / \mathbf{GL}(n, \mathbf{Z})\) .

\subsection{Another interpretation of the topology of the space of closed subgroups}

Let F be the set of closed subsets of G; one defines a Hausdorff uniform structure on F in the following way: for every compact subset K of G and every neighborhood V of e in G, let \(\mathrm{P}(\mathrm{K},\mathrm{V})\) be the set of pairs \((\mathrm{X},\mathrm{Y})\) of elements of F such that both

\[
\mathrm{X} \cap \mathrm{K} \subset \mathrm{VY} \quad \text { and } \quad \mathrm{Y} \cap \mathrm{K} \subset \mathrm{VX}.\tag{9}
\]

Let us show that the set of \(\mathrm{P}(\mathbf{K},\mathbf{V})\) is a fundamental system of entourages for a Hausdorff uniform structure \(\mathcal{U}\) on \(\mathfrak{F}\) . The axioms \((\mathrm{U}_{\mathrm{I}}^{\prime})\) and \((\mathrm{U}_{\mathrm{II}}^{\prime})\) of GT, II, §1, No.~1 are obviously satisfied; moreover, the relations \(\mathbf{K} \subset \mathbf{K}'\) and \(\mathbf{V}' \subset \mathbf{V}\) imply \(\mathrm{P}(\mathbf{K}',\mathbf{V}') \subset \mathrm{P}(\mathbf{K},\mathbf{V})\) ; to verify \((\mathrm{U}_{\mathrm{III}}^{\prime})\) , one can therefore limit oneself to the case that \(\mathbf{V}\) is a symmetric compact neighborhood of \(e\) , so that VK is compact. Suppose that \((\mathbf{X},\mathbf{Y}) \in \mathrm{P}(\mathrm{VK},\mathbf{V})\) and \((\mathbf{Y},\mathbf{Z}) \in \mathrm{P}(\mathrm{VK},\mathbf{V})\) ; then \(\mathbf{X} \cap \mathbf{K} \subset \mathbf{X} \cap \mathrm{VK} \subset \mathrm{VY}\) , and if \(y \in \mathbf{Y}\) is such that \(vy \in \mathbf{K}\) for some \(v \in \mathbf{V}\) , then necessarily \(y \in \mathrm{VK}\) , therefore

\[
\mathrm{X} \cap \mathrm{K} \subset \mathrm{V} (\mathrm{Y} \cap \mathrm{VK});
\]

on the other hand, \(Y \cap VK \subset VZ\) , whence \(X \cap K \subset V^{2}Z\) , and one shows similarly that \(Z \cap K \subset V^{2}X\) , which proves \((U_{III}^{\prime})\) . Finally, if X,Y are two distinct elements of F, there exists for example a point \(a \in X\) such that \(a \notin Y\) , hence a symmetric compact neighborhood V of e such that \(Va \cap Y = \varnothing\) , that is, \(a \notin VY\) ; a fortiori \((X,Y) \notin P(Va,V)\) , which completes the proof of our assertion.

This established, let us consider on the set \(\Sigma\) of closed subgroups of \(G\) the topology \(\mathcal{T}\) induced by the topology of the uniform space \(\mathfrak{F}\) just defined. We shall see that this topology is identical to the topology defined in No.~3. It will suffice to prove that the mapping \(\alpha \mapsto H_{\alpha}\) of \(\Gamma\) into \(\Sigma\) is continuous when \(\Sigma\) is equipped with the topology \(\mathcal{T}\) : for, the same will then be true of the restriction of this mapping to \(\Gamma_{\varphi}\) (with notations as in No.~1, Prop. 2), which is bijective; but since \(\Gamma_{\varphi}\) is compact and the topology \(\mathcal{T}\) is separated, the mapping \(\alpha \mapsto H_{\alpha}\) of \(\Gamma_{\varphi}\) into \(\Sigma\) will then be a homeomorphism.

Thus let \(\alpha_0\) be a point of \(\Gamma\) and let \(\Phi\) be a filter on \(\Gamma\) that converges to \(\alpha_0\) ; we are to show that, with respect to \(\Phi\) , \(\mathrm{H}_{\alpha}\) tends to \(\mathrm{H}_{\alpha_0}\) for the topology \(\mathcal{T}\) . Let \(\mathbf{K}\) be a compact subset of \(\mathbf{G}\) , \(\mathbf{V}\) a symmetric compact neighborhood of \(e\) in \(\mathbf{G}\) ; for every \(x \in \mathrm{H}_{\alpha_0} \cap \mathrm{K}\) , there exists a set \(\mathrm{M}(x) \in \Phi\) such that for every \(\alpha \in \mathrm{M}(x)\) , one has \(\mathrm{V}x \cap \mathrm{H}_{\alpha} \neq \emptyset\) (No.~1, Lemma 2), whence \(\mathrm{V}x \subset \mathrm{V}^2\mathrm{H}_{\alpha}\) ; on covering \(\mathrm{H}_{\alpha_0} \cap \mathrm{K}\) by a finite number of sets \(\mathrm{V}x_i\) , one sees that if \(\mathrm{M} = \bigcap \mathrm{M}(x_i)\) , then \(\mathrm{H}_{\alpha_0} \cap \mathrm{K} \subset \mathrm{V}^2\mathrm{H}_{\alpha}\) for every \(\alpha \in \mathrm{M}\) .

Conversely, suppose that there existed an open neighborhood U of e in G such that, for every set \(L \in \Phi\) , there is at least one \(\alpha \in L\) for which \(H_{\alpha} \cap K \not\subset UH_{\alpha_{0}}\) ; if \(\omega(L)\) is the set of \(\alpha \in L\) having this property, the \(\omega(L)\) would form a base of a filter \(\Phi'\) on \(\Gamma\) finer than \(\Phi\) , and, for every \(\alpha\) belonging to the union E of the \(\omega(L)\) for \(L \in \Phi\) , there would exist a \(t_{\alpha} \in H_{\alpha} \cap K\) not belonging to \(UH_{\alpha_0}\) ; for \(\alpha \notin E\) , take for \(t_{\alpha}\) any point of \(H_{\alpha}\) . Since \(K \cap \mathbb{C}(UH_{\alpha_0})\) is compact, there would exist a cluster point \(s\) of \(\alpha \mapsto t_{\alpha}\) with respect to \(\Phi'\) , belonging to \(K \cap \mathbb{C}(UH_{\alpha_0})\) ; but since \(\Phi'\) converges to \(\alpha_0\) in \(\Gamma\) , this contradicts Lemma 3 of No.~1.

\specialchapter{Exercises}

\exercisesection{}

\begin{enumerate}
\def\labelenumi{\arabic{enumi})}
\item
  Let \(\Gamma\) be a proper \(^{1}\) closed convex cone in \(\mathbf{R}^{n}\) . Show that the mapping \((x,y)\mapsto x+y\) of \(\Gamma\times\Gamma\) into \(\Gamma\) is proper. Deduce from this that any two measures on \(\Gamma\) are convolvable for the mapping \((x,y)\mapsto x+y\) .
\item
  Let G be a locally compact group and \(\Gamma\) the compact space obtained by adjoining to G a point at infinity \(\omega\) . One extends the law of composition of G to \(\Gamma\) by setting \(x\omega = \omega x = \omega\) for every \(x \in \Gamma\) . To every measure \(\mu\) on \(\Gamma\) there corresponds, on the one hand a bounded measure \(\mu_{1}\) on G, on the other a complex number \(\mu(\{\omega\})\) . Show that if \(*(\text{resp.}\widehat{*})\) denotes the convolution defined by the multiplication in G (resp. \(\Gamma\) ), then \((\mu\widehat{*}\nu)_{1} = \mu_{1}*\nu_{1}\) and
\end{enumerate}

\[
(\mu \widehat {*} \nu) (\omega) = \mu (\omega) \nu_ {1} (G) + \nu (\omega) \mu_ {1} (G) + \mu (\omega) \nu (\omega)
\]

for any two measures \(\mu\) and \(\nu\) on \(\Gamma\) .

\exercisesection{}

\begin{enumerate}
\def\labelenumi{\arabic{enumi})}
\tightlist
\item
  Let \((\mathbf{G}_{\iota})_{\iota \in \mathbf{I}}\) be a family of locally compact groups, all but finitely many of them compact. Let \(U_{\iota}\) be a continuous linear representation of \(\mathbf{G}_{\iota}\) in a locally convex space \(\mathbf{E}_{\iota}\) . For every \(s = (s_{\iota}) \in \mathbf{G} = \prod_{\iota} \mathbf{G}_{\iota}\) , let \(U(s)\) be the endomorphism \((x_{\iota}) \mapsto (U_{\iota}(s)x_{\iota})\) of
\end{enumerate}

\(\mathbf{E} = \prod_{\iota}\mathbf{E}_{\iota}\) . Show that \(U\) is a continuous linear representation of \(\mathbf{G}\) in \(\mathbf{E}\) . Let \(\mathbf{E}'\) be the topological direct sum of the \(E_{\iota}\) . Let \(V(s)\) be the restriction of \(U(s)\) to \(E'\) . Show that V is a continuous linear representation of G in \(E'\) .

\begin{enumerate}
\def\labelenumi{\arabic{enumi})}
\setcounter{enumi}{1}
\tightlist
\item
  Let \(U_{1}\) (resp. \(U_{2}\) ) be a continuous linear representation of a locally compact group \(G\) (resp. \(H\) ) in a locally convex space \(E\) (resp. \(F\) ). For \(u \in \mathcal{L}(E; F)\) , \(x \in G\) , \(y \in H\) , set
\end{enumerate}

\[
V (x, y) \cdot u = U _ {2} (y) \circ u \circ U _ {1} (x).
\]

Show that the mapping \((x,y)\mapsto V(x,y)\) is a continuous linear representation of the group \(G^{0}\times H\) in the space \(\mathcal{L}(\mathrm{E};\mathrm{F})\) equipped with the topology of compact convergence. (Use Prop. 9 of TVS, III, §4, No.~4 and the fact that, for K compact in G, \(U_{1}(K)\) is equicontinuous.)

¶ 3) Let G be a locally compact group, U a continuous linear representation of G in a locally convex space E, and E' the dual of E equipped with the strong topology.

\begin{enumerate}
\def\labelenumi{\alph{enumi})}
\item
  Show that for every compact subset K of G, \({}^{t}U(K)\) is equicontinuous.
\item
  Let \(\mathbf{F}\) be the set of \(a' \in \mathbf{E}'\) such that the mapping \(s \mapsto {}^t U(s)a'\) of \(\mathbf{G}\) into \(\mathbf{E}'\) is continuous. Show that \(\mathbf{F}\) is a closed linear subspace of \(\mathbf{E}'\) stable for \({}^t U(\mathbf{G})\) and that the representation deduced by restricting to \(\mathbf{F}\) the contragredient representation of \(U\) is continuous.
\item
  Assume that E is quasi-complete. Let \(\alpha\) be a left Haar measure on G. Show that \(f \mapsto U(f \cdot \alpha)\) is a continuous mapping of \(\mathcal{K}(\mathrm{G})\) into \(\mathcal{L}(\mathrm{E};\mathrm{E})\) equipped with the topology of bounded convergence. (Use Prop. 17 of Ch. VI, §1, No.~7.) Show that F is weakly dense in \(E'\) . (Prove that \({}^{t}U(f)a' \in F\) for all \(a' \in E'\) and \(f \in \mathcal{K}(\mathrm{G})\) , then use Cor. 3 of Lemma 4.) Deduce from this that if E is semi-reflexive, then the contragredient representation of U in \(E'\) equipped with the strong topology is continuous.
\item
  Show that if U is taken to be the left regular representation of G in \(\mathrm{L}^{1}(\mathrm{G},\alpha)\) ( \(\alpha\) still being a left Haar measure on G), then F is the subspace of \(\mathrm{E}' = \mathrm{L}^{\infty}(\mathrm{G},\alpha)\) formed by the uniformly continuous functions.
\end{enumerate}

\begin{enumerate}
\def\labelenumi{\arabic{enumi})}
\setcounter{enumi}{3}
\item
  Let H be a Hilbert space. A continuous representation U of G in H is said to be unitary if the endomorphisms \(U(s)\) are unitary for all \(s \in G\) . For every \(\mu \in \mathcal{M}(G)\) , let \(\mu^{*}\) be the conjugate measure of \(\check{\mu}\) . Show that if \(\mu \in \mathcal{M}^{1}(G)\) then \(U(\mu^{*}) = U(\mu)^{*}\) .
\item
  Let \(\mathbf{G}\) be a locally compact group, \(\mathbf{H}\) a closed subgroup of \(\mathbf{G}\) , and \(U\) a continuous linear representation of \(\mathbf{H}\) in a locally convex space \(\mathbf{E}\) . Let \(\mathbf{K}\) be a compact subset of \(\mathbf{G}\) . Let \(\mathcal{H}^U(\mathbf{K})\) be the space of continuous functions on \(\mathbf{G}\) , with values in \(\mathbf{E}\) , with support contained in \(\mathrm{KH}\) , and satisfying \(f(xh) = U(h)^{-1}f(x)\) ( \(x \in \mathbf{G}, h \in \mathbf{H}\) ). Let \(\mathcal{H}^U\) be the union of the \(\mathcal{H}^U(\mathbf{K})\) , equipped with the direct limit topology of the topologies of uniform convergence in \(\mathbf{K}\) on each of the spaces \(\mathcal{H}^U(\mathbf{K})\) . For \(f \in \mathcal{H}^U\) and \(s \in \mathbf{G}\) , define \(V(s)f \in \mathcal{H}^U\) by
\end{enumerate}

\[
(V (s) f) (t) = f \left(s ^ {- 1} t\right).
\]

Show that \(V\) is a continuous linear representation of \(\mathbf{G}\) in \(\mathcal{K}^U\) .

\begin{enumerate}
\def\labelenumi{\arabic{enumi})}
\setcounter{enumi}{5}
\tightlist
\item
  Let \(\mathbf{G}\) be a locally compact group, \(\beta\) a relatively invariant, nonzero positive measure on \(\mathbf{G}\) , \(\chi\) and \(\chi'\) its left and right multipliers. For \(f \in L_{\mathbf{C}}^{p}(\mathbf{G}, \beta)\) and \(s \in \mathbf{G}\) , set
\end{enumerate}

\[
\begin{array}{c} (U (s) f) (x) = \chi (s) ^ {- 1 / p} f (s ^ {- 1} x) \\ (V (s) f) (x) = \chi^ {\prime} (s) ^ {- 1 / p} f (x s) \\ (S f) (x) = (\chi \chi^ {\prime}) (x) ^ {- 1 / p} \overline {{f (x ^ {- 1})}}. \end{array}
\]

Then \(U\) and \(V\) are linear representations of \(\mathbf{G}\) , and

\[
\begin{array}{c} S ^ {2} = 1, \| U (s) \| = \| V (s) \| = 1 \\ U (s) V (t) = V (t) U (s), S U (s) S = V (s) \end{array}
\]

for all \(s, t\) in \(\mathbf{G}\) .

\begin{enumerate}
\def\labelenumi{\arabic{enumi})}
\setcounter{enumi}{6}
\item
  Let E be a Hilbert space having an orthonormal basis \((e_{s})_{s\in\mathbf{R}}\) equipotent to R. For every \(s\in R\) , denote by \(U(s)\) the isometry of E such that \(U(s)\cdot e_{t}=e_{s+t}\) for all \(t\in R\) ; the linear representation \(s\mapsto U(s)\) of R in E is not continuous, even though the set of \(U(s)\) is equicontinuous.
\item
  Let \(G\) be a locally compact abelian group, \(\mu\) a Haar measure on \(G\) , \(f\) a \(\mu\) -measurable finite numerical function on \(G\) . Assume that for every \(s \in G\) , the numerical function
\end{enumerate}

\[
x \mapsto f (s x) - f (x)
\]

is continuous on G. Show that f is then continuous. (Argue by contradiction; assuming that the function f is not continuous at a point \(x_{0} \in G\) , show first that there exists on G a filter F with limit e such that

\[
\lim _ {\mathfrak {F}, s} | f (s x _ {0}) | = + \infty ,
\]

and deduce from this that also \(\lim_{\mathfrak{F},s}|f(sx)| = +\infty\) for every \(x \in G\) . If \(g = |f| / (1 + |f|)\) , deduce from the last result a contradiction to the fact that for every compact subset \(K \subset G\) ,

\[
\lim _ {\mathfrak {F}, s} \int_ {\mathrm{K}} | g (s x) - g (x) | d \mu (x) = 0.
\]

\begin{enumerate}
\def\labelenumi{\arabic{enumi})}
\setcounter{enumi}{8}
\item
  Let \(G\) be a locally compact group, \(\mu\) a left Haar measure on \(G\) , \(f\) a \(\mu\) -integrable function. Let \(\mathfrak{B}\) be a filter base formed of \(\mu\) -integrable sets of measure \(> 0\) , having \(e\) as limit. For every \(B \in \mathfrak{B}\) , set \(f_{\mathrm{B}}(t) = \frac{1}{\mu(\mathrm{B})} \int_{\mathrm{B}} f(st) d\mu(s)\) . Show that for every integrable subset \(A\) of \(G\) , \(\lim_{\mathfrak{B}} \int_{A} f_{\mathrm{B}}(t) d\mu(t) = \int_{A} f(t) d\mu(t)\) .
\item
  Let \(\mathbf{G}\) be a locally compact group, \(\mathbf{E}\) a Hausdorff locally convex space, \(\mathbf{E}'\) its dual, and \(U\) a linear representation of \(\mathbf{G}\) in \(\mathbf{E}\) , continuous for the weakened topology \(\sigma(\mathbf{E},\mathbf{E}')\) on \(\mathbf{E}\) . Assume that \(\mathbf{E}\) is quasi-complete for \(\sigma(\mathbf{E},\mathbf{E}')\) , so that \(U(\mu)\) is defined for every measure \(\mu \in \mathcal{C}'(\mathbf{G})\) . Show that the bilinear mapping \((\mu,x) \mapsto U(\mu) \cdot x\) is hypocontinuous relative to the equicontinuous subsets of \(\mathcal{C}'(\mathbf{G})\) .
\end{enumerate}

\exercisesection{}

\begin{enumerate}
\def\labelenumi{\arabic{enumi})}
\tightlist
\item
  In \(\mathbf{R}^2\) , let \(\lambda\) and \(\mu\) be the positive measures
\end{enumerate}

\[
f \mapsto \int_ {0} ^ {+ \infty} f (x, 0) d x, \quad f \mapsto \int_ {0} ^ {+ \infty} f (- x, x) d x \quad \left(f \in \mathcal {K} (\mathbf {R} ^ {2})\right).
\]

Show that \(\lambda\) and \(\mu\) are convolvable. Let \(u\) be the homomorphism \((x,y) \mapsto x\) of \(\mathbf{R}^2\) onto \(\mathbf{R}\) . Show that \(u\) is \(\lambda\) -proper and \(\mu\) -proper, but that \(u(\lambda)\) and \(u(\mu)\) are not convolvable.

\begin{enumerate}
\def\labelenumi{\arabic{enumi})}
\setcounter{enumi}{1}
\item
  Let X be a locally compact space in which a locally compact group G operates on the left continuously. Let E be a linear subspace of \(\mathcal{M}(X)\) , stable for the \(\boldsymbol{\gamma}(s)\) ( \(s \in G\) ), equipped with a quasi-complete locally convex topology finer than the topology of compact convergence in \(\mathcal{K}(X)\) . For every \(s \in G\) , let \(\boldsymbol{\gamma}_{\mathrm{E}}(s)\) be the restriction of \(\boldsymbol{\gamma}(s)\) to E. Assume that \(\mu \in E\) implies \(|\mu| \in E\) , and that the representation \(\gamma_{E}\) of G in E is equicontinuous. Then, if \(\mu \in \mathbf{E}\) and \(\nu \in \mathcal{M}^1 (\mathbf{G})\) , \(\nu\) and \(\mu\) are convolvable and \(\nu * \mu = \gamma_{\mathbf{E}}(\nu)\mu \in \mathbf{E}\) . (Make use notably of Prop. 17 of Ch. VI, §1, No.~7.)
\item
  For \(x = (x_{1}, \ldots, x_{n}) \in \mathbf{R}^{n}\) one sets \(|x|^{2} = x_{1}^{2} + \cdots + x_{n}^{2}\) . Let \(M_{1}\) be the set of measures \(\mu\) on \(R^{n}\) for which there exists a real number k such that the function \((1 + |x|^{2})^{k}\) is \(\mu\) -integrable. Let \(M_{2}\) be the set of measures \(\nu\) on \(R^{n}\) for which the function \((1 + |x|^{2})^{k}\) is \(\nu\) -integrable for every k. Show that if \(\mu \in M_{1}\) and \(\nu \in M_{2}\) , then \(\mu\) and \(\nu\) are convolvable, and \(\mu * \nu \in M_{1}\) ; if \(\mu \in M_{2}\) and \(\nu \in M_{2}\) , then \(\mu * \nu \in M_{2}\) . (One first shows that if \(u \geqslant 0\) , \(v \geqslant 0\) , then
\end{enumerate}

\[
(1 + u ^ {2}) (1 + v ^ {2}) \geqslant \frac {1}{3} (1 + (u + v) ^ {2});
\]

from this, one deduces that for any \(x, y\) in \(\mathbf{R}^n\) ,

\[
1 + | x | ^ {2} \leqslant 3 (1 + | y | ^ {2}) (1 + | x + y | ^ {2}).
\]

Then let \(\mu \in \mathcal{M}_1\) , \(\nu \in \mathcal{M}_2\) with \(\mu \geqslant 0\) , \(\nu \geqslant 0\) . Let \(f\) be a continuous function \(\geqslant 0\) on \(\mathbf{R}^n\) . There exists a \(k\) such that \(\mu = (1 + |x|^2)^k \cdot \mu_1\) with \(\mu_1\) bounded; let \(\nu_1 = (1 + |x|^2)^k \cdot \nu\) , which is bounded. Then

\[
\int^ {*} f (x + y) d \mu (x) d \nu (y) \leqslant 3 ^ {k} \int^ {*} (1 + | x | ^ {2}) ^ {k} f (x) d (\mu_ {1} * \nu_ {1}) (x).
\]

Whence the convolvability of \(\mu\) and \(\nu\) and the fact that \(\mu * \nu \in \mathcal{M}_1\) . Argue in an analogous manner for \(\mu, \nu\) in \(\mathcal{M}_2\) .)

\begin{enumerate}
\def\labelenumi{\arabic{enumi})}
\setcounter{enumi}{3}
\item
  Let \(\mu\) be Lebesgue measure on \(\mathbf{R}\) , and \(\nu\) Lebesgue measure on \([0, +\infty[\) . Let \(x_1, x_2\) be in \(\mathbf{R}\) . Show that the convolution product \(((\varepsilon_{x_1} - \varepsilon_{x_2}) * \nu) * \mu\) is defined, but that \(\mu\) and \(\nu\) are not convolvable. Show that the convolution products \(\nu * ((\varepsilon_{x_1} - \varepsilon_{x_2}) * \mu)\) and \((\nu * (\varepsilon_{x_1} - \varepsilon_{x_2})) * \mu\) are defined, but are distinct for \(x_1 \neq x_2\) .
\item
  Let G be a compact group, and \(\mu\) a positive measure on G, with support G, such that \(\mu*\mu=\mu\) . Show that \(\mu\) is the normalized Haar measure of G. (First show that \(\|\mu\|=1\) . Then, if \(\mu\) is not the Haar measure of G, there exists a function \(f\in\mathcal{K}_{+}(G)\) such that
\end{enumerate}

\[
\int f (s) d \mu (s) \geqslant \int f (s t) d \mu (s)
\]

for all \(t \in \mathbf{G}\) , and such that \(\int f(s) d\mu(s) > \int f(st) d\mu(s)\) for some \(t\) . Show that one then has \((\mu * \mu)(f) > \mu(f)\) .

\begin{enumerate}
\def\labelenumi{\arabic{enumi})}
\setcounter{enumi}{5}
\tightlist
\item
  \begin{enumerate}
  \def\labelenumii{\alph{enumii})}
  \tightlist
  \item
    Let \(I = [0,1]\) . Let \(f\) be a continuous function \(\geqslant 0\) on \(\mathbf{R}\) , with support contained in \([-1,0]\) . Show that the set of functions \(\pmb{\gamma}(s)f|I\) ( \(s \in I\) ) has infinite rank in \(\mathcal{K}(I)\) .
  \end{enumerate}
\end{enumerate}

\begin{enumerate}
\def\labelenumi{\alph{enumi})}
\setcounter{enumi}{1}
\item
  Let \(f_1, \ldots, f_n\) be in \(\mathcal{K}(\mathbf{R})\) . Let \(M\) be the set of measures \(\mu \in \mathcal{M}(I)\) such that \(\mu(f_1) = \ldots = \mu(f_n) = 0\) . Show that the set of functions \(y \mapsto \int f(x + y) d\mu(x)\) ( \(y \in I\) ), where \(\mu\) runs over \(M\) , has infinite rank in \(\mathcal{K}(I)\) . (Make use of \(a\) ).
\item
  Let \(g_1, \ldots, g_p\) be in \(\mathcal{K}(\mathbf{R})\) . Deduce from \(b\) that there exist \(\mu \in \mathbb{M}\) and \(\nu \in \mathcal{M}(\mathrm{I})\) such that \(\nu(g_1) = \ldots = \nu(g_p) = 0\) , \((\nu * \mu)(f) \neq 0\) .
\item
  Deduce from c) that the mapping \((\mu,\nu)\mapsto\nu*\mu\) of \(\mathcal{M}(\mathbf{T})\times\mathcal{M}(\mathbf{T})\) into \(\mathcal{M}(\mathbf{T})\) is not vaguely continuous.
\end{enumerate}

\begin{enumerate}
\def\labelenumi{\arabic{enumi})}
\setcounter{enumi}{6}
\tightlist
\item
  Let \(G\) be a locally compact group.
\end{enumerate}

\begin{enumerate}
\def\labelenumi{\alph{enumi})}
\item
  Show that if a positive measure \(\mu \in \mathcal{C}'(\mathrm{G})\) admits an inverse in \(\mathcal{C}'(\mathrm{G})\) that is a positive measure, then \(\mu\) is necessarily a point measure.
\item
  Taking G to be the finite group \(\mathbf{Z}/2\mathbf{Z}\) , give an example of a non point measure \(\mu \in \mathcal{M}(\mathbf{G})\) that is invertible in \(\mathcal{M}(\mathbf{G})\) .
\end{enumerate}

\begin{enumerate}
\def\labelenumi{\arabic{enumi})}
\setcounter{enumi}{7}
\item
  Let G be a locally compact group; denote by \(\mathcal{C}_{+}^{\prime}(G)\) the set of positive measures on G with compact support. Show that the mapping \((\mu, \nu) \mapsto \mu * \nu\) of \(\mathcal{M}_{+}(G) \times \mathcal{C}_{+}^{\prime}(G)\) into \(\mathcal{M}_{+}(G)\) is continuous when \(\mathcal{C}'(G)\) is equipped with the weak topology \(\sigma(\mathcal{C}'(G), \mathcal{C}(G))\) , and \(\mathcal{M}(G)\) with the vague topology \(\sigma(\mathcal{M}(G), \mathcal{K}(G))\) . (Make use of Exer. 10 of Ch. III, §1 and the fact that G is paracompact.)
\item
  \begin{enumerate}
  \def\labelenumii{\alph{enumii})}
  \tightlist
  \item
    Let G be a locally compact group, B a bounded subset of \(\mathcal{M}(\mathbf{G})\) , C an equicontinuous subset of \(\mathcal{C}'(\mathbf{G})\) ; show that if \(\mathcal{C}'(\mathbf{G})\) is equipped with the weak topology \(\sigma(\mathcal{C}'(\mathbf{G}), \mathcal{C}(\mathbf{G}))\) , and \(\mathcal{M}(\mathbf{G})\) with the vague topology \(\sigma(\mathcal{M}(\mathbf{G}), \mathcal{K}(\mathbf{G}))\) , then the mapping \((\mu, \nu) \mapsto \mu * \nu\) of B × C into \(\mathcal{M}(\mathbf{G})\) is continuous. (Observe that all of the measures \(\nu \in \mathbf{C}\) have their support in a common compact set, and make use of Ch. III, §4, No.~3, remarks following Prop. 6.)
  \end{enumerate}
\end{enumerate}

\begin{enumerate}
\def\labelenumi{\alph{enumi})}
\setcounter{enumi}{1}
\tightlist
\item
  If, in \(\mathcal{M}^1 (\mathbf{R})\) , one sets \(\mu_n = \varepsilon_n\) , \(\nu_{n} = \varepsilon_{-n}\) , the sequences \((\mu_n)\) and \((\nu_{n})\) tend vaguely to 0 for the weak topology \(\sigma (\mathcal{M}^1 (\mathbf{R}),\overline{\mathcal{K}(\mathbf{R})})\) , but the sequence \((\mu_n*\nu_n)\) does not tend to 0 for the vague topology.
\end{enumerate}

¶ 10) For a locally compact space T, one denotes by \(\mathcal{C}^{\infty}(T)\) the Banach space of continuous bounded numerical functions on T. A subset H of \(\mathcal{M}^{1}(T)\) is said to be cramped if, for every \(\varepsilon > 0\) , there exists a compact subset K of T such that \(|\mu|(T - K) \leqslant \varepsilon\) for all \(\mu \in H\) .

\begin{enumerate}
\def\labelenumi{\alph{enumi})}
\item
  Show that if a subset H of \(\mathcal{M}^{1}(G)\) is cramped and is bounded for the topology defined by the norm of \(\mathcal{M}^{1}(T)\) , then it is relatively compact for the topology \(\sigma(\mathcal{M}^{1}(T),\mathcal{C}^{\infty}(T))\) (observe that H is relatively compact for \(\sigma(\mathcal{M}^{1}(T),\mathcal{K}(T))\) ).
\item
  Assuming in addition that T is paracompact, show that, conversely, if H is a subset of \(\mathcal{M}^1 (\mathrm{T})\) relatively compact for \(\sigma (\mathcal{M}^1 (\mathrm{T}),\mathcal{C}^\infty (\mathrm{T}))\) , then H is bounded for the norm of \(\mathcal{M}^1 (\mathrm{T})\) and is cramped. (Consider first the case that \(\mathbf{T} = \mathbf{N}\) , applying Exer. 15 of Ch. V, §5. Then contemplate the case that T is countable at infinity, the union of a sequence \((\mathrm{U}_n)\) of relatively compact open sets such that \(\overline{\mathrm{U}}_n\subset \mathrm{U}_{n + 1}\) . Arguing by contradiction, show that one can reduce to the case that there would exist for each \(n\) a continuous numerical function \(f_{n}\) defined on T, with support contained in \(\mathrm{U}_{n + 1} - \overline{\mathrm{U}}_n\) , such that \(\| f_n\| \leqslant 1\) , and a sequence \((\mu_n)\) of measures belonging to H and for which \(\mu_{n}(f_{n})\geqslant a > 0\) for all \(n\) . Consider then the continuous mapping \(u:\mathrm{L}^{1}(\mathbf{N})\to \mathcal{C}^{\infty}(\mathrm{T})\)
\end{enumerate}

such that \(u((\xi_n)) = \sum_{n=0}^{\infty} \xi_n f_n\) and obtain a contradiction to what was proved for \(T = N\) .

by considering the transpose of \(u\) . Finally, when \(T\) is an arbitrary paracompact locally compact space, \(T\) is the topological sum of a family \((T_{\alpha})\) of locally compact spaces that are countable at infinity; for every \(\alpha\) , let \(m_{\alpha} = \sup_{\mu \in H} |\mu|(T_{\alpha})\) ; show, arguing by contradiction and making use of the preceding case, that necessarily \(m_{\alpha} = 0\) except for a countable infinity of indices \(\alpha\) .)

\begin{enumerate}
\def\labelenumi{\alph{enumi})}
\setcounter{enumi}{2}
\tightlist
\item
  Show that the conclusion of \(b\) ) is not valid for the non-paracompact locally compact space defined in Exer. 16 h) of Ch. IV, §4.
\end{enumerate}

¶ 11) Let G be a locally compact group; on \(\mathcal{M}^{1}(G)\) , let us denote by \(\mathcal{T}_{\mathrm{I}}\) the topology \(\sigma(\mathcal{M}^{1}(G), \overline{\mathcal{K}(G)})\) , by \(\mathcal{T}_{\mathrm{III}}\) the topology \(\sigma(\mathcal{M}^{1}(G), \mathcal{C}^{\infty}(G))\) , and let us write \(\mathcal{M}_{\mathrm{I}}\) , \(\mathcal{M}_{\mathrm{III}}\) for the space \(\mathcal{M}^{1}(G)\) equipped respectively with \(\mathcal{T}_{\mathrm{I}}\) and \(\mathcal{T}_{\mathrm{III}}\) .

\begin{enumerate}
\def\labelenumi{\alph{enumi})}
\item
  Let A be a bounded subset of \(\mathcal{M}_{\mathrm{I}}\) , and B a relatively compact subset of \(\mathcal{M}_{\mathrm{III}}\) . Show that the restriction to \(\mathrm{A} \times \mathrm{B}\) of the mapping \((\mu, \nu) \mapsto \mu * \nu\) of \(\mathcal{M}_{\mathrm{I}} \times \mathcal{M}_{\mathrm{III}}\) into \(\mathcal{M}_{\mathrm{I}}\) is continuous (use Exer. 10 to reduce to evaluating an integral \(\iint f(st) d\mu(s) d\nu(t)\) when \(f(st) = \sum_{i} u_{i}(s)v_{i}(t)\) with \(u_{i}\) and \(v_{i}\) in \(\mathcal{K}(\mathrm{G})\) ).
\item
  Give an example where G is compact (hence \(T_{I} = T_{III}\) ) showing that \((\mu, \nu) \mapsto \mu * \nu\) , regarded as a mapping of \(M_{I} \times M_{III}\) into \(M_{I}\) , is not hypocontinuous with respect to the relatively compact subsets of \(M_{I}\) , nor with respect to the relatively compact subsets of \(M_{III}\) (cf.~Exer. 6).
\item
  Let \(A, B\) be two relatively compact subsets of \(\mathcal{M}_{\mathrm{III}}\) . Show that the restriction to \(A \times B\) of the mapping \((\mu, \nu) \mapsto \mu * \nu\) of \(\mathcal{M}_{\mathrm{III}} \times \mathcal{M}_{\mathrm{III}}\) into \(\mathcal{M}_{\mathrm{III}}\) is continuous (same method as in a)).
\item
  Denote by E the subspace of \(\mathbf{R}^{\mathrm{G}}\) formed by the linear combinations of characteristic functions of open subsets of G, by \(\mathcal{T}_{\mathrm{IV}}\) the topology \(\sigma(\mathcal{M}^{1}(\mathrm{G}),\mathrm{E})\) , and by \(\mathcal{M}_{\mathrm{IV}}\) the space \(\mathcal{M}^{1}(\mathrm{G})\) equipped with \(\mathcal{T}_{\mathrm{IV}}\) . Recall that the topologies induced by \(\mathcal{T}_{\mathrm{III}}\) and \(\mathcal{T}_{\mathrm{IV}}\) on a bounded subset of \(\mathcal{M}^{1}(\mathrm{G})\) consisting of positive measures, are in general distinct (Ch. V, §5, Exer. 16 c)). Recall also that the compact subsets of \(\mathcal{M}_{\mathrm{IV}}\) are the same as those of \(\mathcal{M}^{1}(\mathrm{G})\) for the weakened topology \(\sigma(\mathcal{M}^{1}(\mathrm{G}),(\mathcal{M}^{1}(\mathrm{G}))'\) on the Banach space \(\mathcal{M}^{1}(\mathrm{G})\) (Ch. VI, §2, Exer. 12). Show that if A, B are two relatively compact subsets of \(\mathcal{M}_{\mathrm{IV}}\) , the restriction to A × B of the mapping \((\mu,\nu)\mapsto\mu*\nu\) of \(\mathcal{M}_{\mathrm{IV}}\times\mathcal{M}_{\mathrm{IV}}\) into \(\mathcal{M}_{\mathrm{IV}}\) is continuous. (Restrict attention to the case that A and B are compact, and begin by proving that the image of A × B under the preceding mapping is then compact for \(\sigma(\mathcal{M}^{1}(\mathrm{G}),(\mathcal{M}^{1}(\mathrm{G}))'\) ); for this, apply the theorems of Eberlein and Šmulian (TVS, IV, §5, No.~3), and thus reduce to proving that if \((\mu_{n})\) , \((\nu_{n})\) are two sequences that converge to 0 in \(\mathcal{M}_{\mathrm{IV}}\) , then the same is true of the sequence \((\mu_{n}*\nu_{n})\) ; make use of Prop. 12 and Exer. 15 of Ch. V, §5. Finally, to prove that \((\mu,\nu)\mapsto\mu*\nu\) is continuous for \(\mathcal{T}_{\mathrm{IV}}\) , use c) and the fact that \(\mathcal{T}_{\mathrm{III}}\) is coarser than \(\mathcal{T}_{\mathrm{IV}}\) .)
\item
  Take \(\mathbf{G} = \mathbf{R}^2\) ; let \(a, b\) be the vectors of the canonical basis of \(\mathbf{G}\) over \(\mathbf{R}\) , \(\rho_n\) the measure on the interval \(\mathbf{I} = [0, \pi]\) of \(\mathbf{R}\) having as density with respect to Lebesgue measure the function \(\sin(2^n x)\) , and \(\mu_n\) the measure \(\rho_n \otimes \varepsilon_0\) on \(\mathbf{R} \times \mathbf{R}\) ; on the other hand let \(\nu_n = \varepsilon_{b/2^n} - \varepsilon_0\) on \(\mathbf{G}\) . Show that the sequence \((\mu_n)\) tends to 0 in \(\mathcal{M}_{\mathrm{IV}}\) and that the sequence \((\nu_n)\) tends to 0 in \(\mathcal{M}_{\mathrm{III}}\) , but that the sequence \((\mu_n * \nu_n)\) does not tend to 0 in \(\mathcal{M}_{\mathrm{IV}}\) (cf.~Ch. V, §5, Exer. 16).
\item
  Give an example where G is compact and where \((\mu,\nu)\mapsto\mu*\nu\) , regarded as a mapping of \(M_{IV}\times M_{IV}\) into \(M_{IV}\) , is not hypocontinuous relative to the compact subsets of \(M_{IV}\) (same method as in Exer. 6, on observing that for a given \(f\in E\) there exist compact subsets \(H\subset M_{IV}\) such that the set of functions \(t\mapsto\int f(st)d\mu(s)\) , where \(\mu\) runs over H, has finite rank over R).
\end{enumerate}

\begin{enumerate}
\def\labelenumi{\arabic{enumi})}
\setcounter{enumi}{11}
\tightlist
\item
  Let \(G\) be a locally compact group that is not unimodular.
\end{enumerate}

\begin{enumerate}
\def\labelenumi{\alph{enumi})}
\item
  Show that there exists a bounded positive measure \(\mu\) on \(\mathbf{G}\) such that \(\Delta_{\mathbf{G}} \cdot \mu\) is not bounded (take \(\mu\) to be discrete).
\item
  Let \(\mu'\) be a left Haar measure on G. Then \(\mu\) and \(\mu'\) are convolvable (Prop. 5). Show that \(\mu'\) and \(\mu\) are not convolvable.
\end{enumerate}

\begin{enumerate}
\def\labelenumi{\arabic{enumi})}
\setcounter{enumi}{12}
\tightlist
\item
  Let \(r\) be a number such that \(0 < r < 1\) ; for every integer \(n \geqslant 1\) , denote by \(\lambda_{n,r}\) the measure \((\varepsilon_{r^n} + \varepsilon_{-r^n}) / 2\) on \(\mathbf{R}\) , and set \(\mu_{n,r} = \lambda_{1,r} * \lambda_{2,r} * \dots * \lambda_{n,r}\) .
\end{enumerate}

\begin{enumerate}
\def\labelenumi{\alph{enumi})}
\item
  Show that the sequence \((\mu_{n,r})\) converges vaguely to a measure \(\mu_{r}\) on R, with support contained in I = {[}-1,+1{]} (prove that for every interval U of R, the sequence \((\mu_{n,r}(U))\) is convergent).
\item
  Show that for \(r < 1/2\) , the measure \(\mu_r\) is alien to the Lebesgue measure on \(\mathbf{R}\) , but that \(\mu_{1/2}\) is the measure induced on I by Lebesgue measure.
\item
  Let \(\nu_{1/4}\) be the image of \(\mu_{1/4}\) under the homothety \(t \mapsto 2t\) in \(\mathbf{R}\) . Show that \(\mu_{1/4} * \nu_{1/4} = \mu_{1/2}\) , even though \(\mu_{1/4}\) and \(\nu_{1/4}\) are both alien to Lebesgue measure (use Exer. 11 a)).
\end{enumerate}

\exercisesection{}

¶ 1) Let G be a locally compact group, β a left Haar measure on G, A a β-measurable subset of G, and ν a nonzero positive measure on G. Assume that sA is locally ν-negligible for every \(s \in G\) . Show that A is locally β-negligible. (Reduce to the case that A is relatively compact and ν is bounded. Prove that ν and \(\varphi_{A} \cdot \beta\) are convolvable and that \(\nu *^{\beta} \varphi_{A} = 0\) , whence \(0 = \|\nu * \varphi_{A} \cdot \beta\| = \|\nu\| \cdot \|\varphi_{A} \cdot \beta\|\) .) Show, on admitting the continuum hypothesis, that this result may fail to hold if A is not assumed to be β-measurable. (Take \(G = R^{2}\) , take for ν the Haar measure on the subgroup \(R \times \{0\}\) , and apply Exer. 7 c) of Ch. V, §8.)

\begin{enumerate}
\def\labelenumi{\arabic{enumi})}
\setcounter{enumi}{1}
\tightlist
\item
  Let H be the additive group R equipped with the discrete topology. Let G be the locally compact group \(R \times R \times H\) . Let \(\alpha\) be a Haar measure on G, \(\beta\) a Haar measure on \(R \times H\) , and \(\mu = \varepsilon_{0} \otimes \beta \in \mathcal{M}(G)\) . Construct a function \(f \geqslant 0\) on G, locally \(\alpha\) -negligible (hence such that \(\mu\) and f are convolvable) but such that no translate of f is \(\mu\) -measurable. (Imitate the construction of Ch. V, §3, Exer. 4.)
\end{enumerate}

¶ 3) Let G be a locally compact group.

\begin{enumerate}
\def\labelenumi{\alph{enumi})}
\tightlist
\item
  Let \(\mu\) be a nonzero bounded positive measure on \(\mathbf{G}\) such that \(\mu * \mu = \mu\) . Show that the support \(\mathbf{S}\) of \(\mu\) is compact. (Let \(f \in \mathcal{K}_{+}(\mathbf{G})\) with \(f \neq 0\) ; choose a left Haar measure on \(\mathbf{G}\) , with respect to which convolution products will be taken; we have \(\mu * f \in \overline{\mathcal{K}(\mathbf{G})}\) ; let \(x \in \mathbf{S}\) be such that \((\mu * f)(x) = \sup_{y \in \mathbf{S}} (\mu * f)(y)\) ; show that
\end{enumerate}

\((\mu * f)(y) = (\mu * f)(x)\) for every \(y \in S\) .

\begin{enumerate}
\def\labelenumi{\alph{enumi})}
\setcounter{enumi}{1}
\tightlist
\item
  Show that S is a compact subgroup of G and that \(\mu\) is the normalized Haar measure of S. (Make use of a), Exer. 21 of GT, III, §4, and Exer. 5 of §3.)
\end{enumerate}

\begin{enumerate}
\def\labelenumi{\arabic{enumi})}
\setcounter{enumi}{3}
\tightlist
\item
  Let \(G\) be a locally compact group. For every \(t \in \mathbf{R}_+^*\) , let \(\mu_t\) be a nonzero bounded positive measure on \(G\) . Assume that the mapping \(t \mapsto \mu_t\) is continuous for the topology \(\sigma(\mathcal{M}^1(G), \overline{\mathcal{K}(G)})\) , and that \(\mu_{s+t} = \mu_s * \mu_t\) ( \(s, t\) in \(\mathbf{R}_+^*\) ).
\end{enumerate}

\begin{enumerate}
\def\labelenumi{\alph{enumi})}
\item
  Show that there exists a number \(c \in \mathbf{R}\) such that \(\| \mu_t \| = \exp(ct)\) . (Observe that \(t \mapsto \| \mu_t \|\) is lower semi-continuous and that \(\| \mu_{s+t} \| = \| \mu_s \| \cdot \| \mu_t \|\) . Make use of Prop. 18.)
\item
  Suppose that \(c = 0\) . Show that \(\mu_t\) converges weakly as \(t \to 0\) to the normalized Haar measure of a compact subgroup of \(G\) . (Make use of the weak compactness of the unit ball of \(\mathcal{M}^1(G)\) , and show that for every weak cluster point \(\mu\) of \(t \mapsto \mu_t\) with respect to the filter of neighborhoods of 0 in \(R_+^*\) , one has \(\mu_t * \mu = \mu_t\) for every \(t\) , then \(\mu * \mu = \mu\) ; next, apply Exer. 3.)
\end{enumerate}

¶ 5) Let G be a locally compact group countable at infinity, β a left Haar measure on G, \(a \in R_{+}^{*}\) , and \((\nu_{u})_{0 \leqslant u \leqslant a}\) a family of positive measures on G satisfying the following conditions:

\begin{enumerate}
\def\labelenumi{(\roman{enumi})}
\item
  \(\nu_0 = \varepsilon_e\) ; \(\nu_{u + v} = \nu_u * \nu_v\) if \(u + v \leqslant a\) ; \(\nu_u = \stackrel{\vee}{\nu}_u\) ;
\item
  for \(0 < u \leqslant a\) , \(\nu_{u} = f_{u} \cdot \beta\) with \(f_{u} \geqslant 0\) lower semi-continuous;
\item
  \(\nu_{u}\) is a vaguely continuous function of \(u\) .
\end{enumerate}

\begin{enumerate}
\def\labelenumi{\alph{enumi})}
\tightlist
\item
  Let \(f \in \mathcal{K}_{+}(\mathrm{G})\) . Show that, for \(0 < u \leqslant a\) ,
\end{enumerate}

\[
\int \left(f _ {u / 2} * f\right) (z) ^ {2} d \beta (z) = \int f (x) \left(f _ {u} * f\right) (x) d \beta (x).\tag{1}
\]

\begin{enumerate}
\def\labelenumi{\alph{enumi})}
\setcounter{enumi}{1}
\item
  Let \(f \in \mathcal{K}(\mathbf{G})\) . Show that \(f_{u/2} * f\) has \(\beta\) -integrable square for \(0 < u \leqslant a\) , and that (1) again holds.
\item
  Let \(f \in \mathcal{H}(\mathrm{G})\) be such that \(\nu_{a} * f = 0\) . Show that \(f = 0\) . (One has \(\nu_{a/2^n} * f = 0\) for every integer \(n > 0\) by \(b\) ), therefore \(f = 0\) by §3, No.~3, Remark 1).
\item
  Let \(\nu \in \mathcal{C}'(\mathrm{G})\) be such that \(\nu_{a} * \nu = 0\) . Show that \(\nu = 0\) . (Regularize \(\nu\) by functions in \(\mathcal{H}(\mathrm{G})\) and apply \(c\) ).)
\end{enumerate}

\begin{enumerate}
\def\labelenumi{\arabic{enumi})}
\setcounter{enumi}{5}
\item
  Let G be a locally compact group. Show that if \(\mathcal{K}(\mathrm{G})\) is commutative for convolution, then G is abelian. (Show by regularization that \(\mathcal{C}'(\mathrm{G})\) is commutative, and apply this to the measures \(\varepsilon_{s}\) , where \(s \in G\) .)
\item
  Let \(G\) be a locally compact group and \(\beta\) a left Haar measure on \(G\) . Show that the algebra \(L^1(G, \beta)\) has a unity element if and only if \(G\) is discrete. (Suppose \(G\) non-discrete, and let \(f_0 \in L^1(G, \beta)\) . There exists a compact neighborhood \(V\) of \(e\) such that
\end{enumerate}

\[
\int_ {V} | f _ {0} (x) | d \beta (x) <   1.
\]

Let \(\mathbf{U}\) be a symmetric compact neighborhood of \(e\) such that \(\mathbf{U}^2 \subset \mathbf{V}\) . Then, for almost every \(x \in \mathbf{U}\) ,

\[
\left| \left(\varphi_ {\mathrm{U}} * f _ {0}\right) (x) \right| = \int_ {\mathrm{U}} \left| f _ {0} \left(y ^ {- 1} x\right) \right| d \beta (y) \leqslant \int_ {\mathrm{V}} \left| f _ {0} (x) \right| d \beta (x) <   1,
\]

thus \(f_0\) is not a unity element for \(L^1(G, \beta)\) .

\begin{enumerate}
\def\labelenumi{\arabic{enumi})}
\setcounter{enumi}{7}
\tightlist
\item
  Let G be a locally compact group countable at infinity, operating continuously on the left in a Polish locally compact space T. Let \(\nu\) be a positive measure on T that is quasi-invariant under G. Let R be a \(\nu\) -measurable equivalence relation on T, compatible with G. There then exists (Ch. VI, §3, No.~4, Prop. 2) a Polish locally compact space B, and a \(\nu\) -measurable mapping p of T into B, such that \(R\{x,y\}\) is equivalent to \(p(x)=p(y)\) . Let \(\nu'\) be a pseudo-image measure of \(\nu\) under p, and let \(b\mapsto\lambda_{b}\) ( \(b\in B\) ) be a disintegration of \(\nu\) by R. Show that the \(\lambda_{b}\) are, for almost every \(b\in B\) , quasi-invariant under G.
\end{enumerate}

(Let \(\chi\) be a function on \(\mathbf{G} \times \mathbf{T}\) satisfying the conditions of Exer. 13 of Ch. VII, §1. Show that for every \(s \in \mathbf{G}\) , there exists a \(\nu'\) -negligible subset \(\mathbf{N}(s)\) of B such that \(\chi(s^{-1}, \cdot)\) is locally \(\lambda_b\) -integrable for \(b \notin \mathbf{N}(s)\) ; for this, observe that for every \(\psi \in \mathcal{K}(\mathbf{T})\) , the function \(x \mapsto \psi(x)\chi(s^{-1}, x)\) is \(\lambda_b\) -integrable except for \(b\) belonging to a \(\nu'\) -negligible set \(\mathbf{N}(s, \psi)\) , and make use of Lemma 1 of Ch. VI, §3, No.~1. Set \(\lambda_{b,s}' = 0\) if \(b \in \mathbf{N}(s)\) , and \(\lambda_{b,s}' = \chi(s^{-1}, \cdot) \cdot \lambda_b\) if \(b \notin \mathbf{N}(s)\) . Show that the mapping \(b \mapsto \lambda_{b,s}'\) is \(\nu'\) -adequate (use Lemma 3 of Ch. VI, §3, No.~1) and that \(\gamma(s)\beta = \int \lambda_{b,s}' d\nu'(b)\) . Show that on the other hand \(\gamma(s)\beta = \int \gamma(s)\lambda_b d\nu'(b)\) and deduce from this that, for every \(s \in \mathbf{G}\) , one has \(\gamma(s)\lambda_b = \chi(s^{-1}, \cdot) \cdot \lambda_b\) for almost every \(b\) , therefore for almost every \(b\) one has \(\gamma(s)\lambda_b = \chi(s^{-1}, \cdot) \cdot \lambda_b\) for almost every \(s\) . Then use Cor. 2 of Prop. 17 of §4 to infer that, for almost every \(b\) , \(\gamma(s)\lambda_b\) is equivalent to \(\lambda_b\) for all \(s \in \mathbf{G}\) .)

Show that if \(\nu\) is relatively invariant under G with multiplier \(\chi\) , then the \(\lambda_{b}\) are, for almost every b, relatively invariant with multiplier \(\chi\) .

\begin{enumerate}
\def\labelenumi{\arabic{enumi})}
\setcounter{enumi}{8}
\item
  Let \(G\) be a locally compact group, \(\beta\) a left Haar measure on \(G\) , \(A\) and \(B\) two \(\beta\) -integrable sets such that \(\beta(A) \leqslant \beta(B)\) and \(\beta^{*}(A^{-1}) < +\infty\) . Show that there exist
\item
  disjoint sets \(N, K_{1}, K_{2}, \ldots\) covering \(A\) , with \(N \beta\) -negligible and the \(K_{n}\) compact;
\item
  disjoint sets \(\mathbf{N}^{\prime}\) , \(\mathbf{K}_1^{\prime}\) , \(\mathbf{K}_2^{\prime}\) , \ldots{} covering B, with the \(\mathbf{K}_n^{\prime}\) compact;\\
\item
  elements \(s_n \in G\) such that \(\mathbf{K}_n^{\prime} = s_n\mathbf{K}_n\) .
\end{enumerate}

(Using the fact that \(\beta (x\mathbf{A}\cap \mathbf{B})\) depends continuously on \(x\) and that

\[
\int \beta (x \mathrm{A} \cap \mathrm{B}) d \beta (x) = \beta (\mathrm{A} ^ {- 1}) \beta (\mathrm{B}),
\]

show that if \(\beta (\mathbf{A})\neq 0\) then there exists an \(x\in \mathbf{G}\) such that \(\beta (x\mathbf{A}\cap \mathbf{B})\neq 0\) .)

\begin{enumerate}
\def\labelenumi{\arabic{enumi})}
\setcounter{enumi}{9}
\tightlist
\item
  Let \(\mathbf{G}\) be a locally compact group, \(\beta\) a left Haar measure on \(\mathbf{G}\) . For any two \(\beta\) -integrable sets \(\mathbf{A}, \mathbf{B}\) set \(\rho(\mathbf{A}, \mathbf{B}) = \beta(\mathbf{A} \cup \mathbf{B}) - \beta(\mathbf{A} \cap \mathbf{B})\) .
\end{enumerate}

\begin{enumerate}
\def\labelenumi{\alph{enumi})}
\item
  Let A be a \(\beta\) -integrable set. Show that \(x \mapsto \rho(xA, A)\) is a continuous function.
\item
  Let U be a neighborhood of e. Show that there exist a compact set A and a number \(\varepsilon > 0\) such that \(\rho(x\mathbf{A}, \mathbf{A}) < \varepsilon\) implies \(x \in U\) . (Take for A a neighborhood of e such that \(A \cdot A^{-1} \subset U\) .)
\item
  For a subset C of G to be relatively compact, it is necessary and sufficient that there exist a \(\beta\) -integrable set A and a number \(a (0 < a < 2\beta(A))\) such that \(x \in C\) implies \(\rho(xA, A) \leqslant a\) .
\end{enumerate}

\begin{enumerate}
\def\labelenumi{\arabic{enumi})}
\setcounter{enumi}{10}
\tightlist
\item
  \begin{enumerate}
  \def\labelenumii{\alph{enumii})}
  \tightlist
  \item
    Let G be a locally compact group generated by a compact neighborhood of e. Let \(\varphi\) be a non-surjective continuous endomorphism of G belonging to the closure, for the topology of compact convergence, of the group G of (bicontinuous) automorphisms of G. Then \(\lim_{\psi\in\mathcal{G},\psi\to\varphi}\mod\psi=0\) . (Let K be a compact neighborhood of e that generates G. Let \(\mu\) be a left Haar measure on G. If \(\mu(\varphi(K)) > 0\) , then \(\varphi(K) \cdot \varphi(K)^{-1}\) is a neighborhood of \(e\) in G, therefore \(\varphi(G)\) is an open subgroup of G; for \(\psi \in \mathcal{G}\) sufficiently near \(\varphi\) , one has \(\psi(K) \subset \varphi(G)\) , therefore \(\psi(G) \neq G\) , which is absurd. Thus \(\mu(\varphi(K)) = 0\) . As \(\psi \in \mathcal{G}\) tends to \(\varphi\) , \(\mu(\psi(K))\) tends to 0.)
  \end{enumerate}
\end{enumerate}

\begin{enumerate}
\def\labelenumi{\alph{enumi})}
\setcounter{enumi}{1}
\tightlist
\item
  Let G be a free abelian group, a direct sum \(G_{1} \oplus G_{2} \oplus \cdots\) , where each \(G_{i}\) is isomorphic to Z. Consider G as being discrete. Let \(\varphi\) be the non-surjective endomorphism \((x_{1}, x_{2}, \ldots) \mapsto (0, x_{1}, x_{2}, \ldots)\) of G. Then \(\varphi\) is the limit, for the topology of compact convergence, of automorphisms of G, and every automorphism of G has modulus 1.
\end{enumerate}

\begin{enumerate}
\def\labelenumi{\arabic{enumi})}
\setcounter{enumi}{11}
\tightlist
\item
  For \(t > 0\) and \(x \in \mathbf{R}\) , let \(F_t(x) = te^{-\pi t^2 x^2}\) . Let \(f \in \overline{\mathcal{H}(\mathbf{R})}\) . Show that \(f * F_t\) tends to \(f\) uniformly as \(t\) tends to \(+\infty\) . (Show that the measures \(F_t(x) dx\) satisfy the conditions of §2, No.~7, Lemma 4, with \(a = 0\) .)
\end{enumerate}

¶ 13) Let G be a locally compact group operating continuously on the left in a Polish locally compact space T. Let β be a left Haar measure on G, ν a quasi-invariant positive measure on T, and χ(s,x) a function \textgreater0 on G × T satisfying the conditions of Exer. 13 of Ch. VII, §1.

For \(f \in \mathrm{L}^p(\mathrm{T}, \nu)\) ( \(1 \leqslant p < +\infty\) ), set

\[
\left(\gamma_ {\chi , p} (s) f\right) (x) = \chi (s ^ {- 1}, x) ^ {1 / p} f (s ^ {- 1} x).
\]

\begin{enumerate}
\def\labelenumi{\alph{enumi})}
\item
  Show that for every \(s \in G\) , \(\pmb{\gamma}_{\chi,p}(s)\) is an isometric endomorphism of \(L^p(T, \nu)\) . Show that the mapping \(s \mapsto \pmb{\gamma}_{\chi,p}(s)\) is a linear representation of \(G\) in \(L^p(T, \nu)\) (argue as in §2, No.~5).
\item
  Let \(f \in \mathcal{L}^p(\mathbf{T}, \nu)\) , and let \(h \in \mathcal{K}(\mathbf{G})\) . Show that the function
\end{enumerate}

\[
(x, s) \mapsto f (s ^ {- 1} x) h (s) \chi (s ^ {- 1}, x) ^ {1 / p}
\]

is \(p\) -th power integrable for \(\beta \otimes \nu\) (begin with the case \(p = 1\) , and use Lemma 1 of §4, No.~1). Show that if \(q\) is the exponent conjugate to \(p\) , and if \(g \in \mathcal{L}^q(\mathrm{T}, \nu)\) , then \(f(s^{-1}x)h(s)g(x)\chi(s^{-1},x)^{1/p}\) is integrable for \(\beta\otimes\nu\) (write \(h\) in the form \(h_1h_2\) with \(h_1,h_2\) in \(\mathcal{K}(\mathrm{G})\) ). Then deduce from the Lebesgue-Fubini theorem that, for \(f\in\mathrm{L}^{p}(\mathrm{T},\nu)\) and \(g\in\mathrm{L}^{q}(\mathrm{T},\nu)\) , the function

\[
s \mapsto \int g (x) (\gamma_ {\chi , p} (s) f) (x) d \nu (x)
\]

is \(\beta\) -measurable.

\begin{enumerate}
\def\labelenumi{\alph{enumi})}
\setcounter{enumi}{2}
\item
  Show, using Cor. 2 of Prop. 18 of §4, No.~6 and Lemma 1 of Ch. VI, §3, No.~1, that the representation \(s \mapsto \gamma_{\chi,p}(s)\) of G in \(L^p(T, \nu)\) is continuous for \(1 \leqslant p < +\infty\) .
\item
  Assume in addition that, for every \(s \in G\) , the function \(x \mapsto \chi(s^{-1}, x)\) is bounded. For \(f \in L^{p}(T, \nu)\) , set
\end{enumerate}

\[
\left(\gamma_ {\chi} (s) f\right) (x) = \chi (s ^ {- 1}, x) f (s ^ {- 1} x).
\]

Show that \(s \mapsto \gamma_{\chi}(s)\) is a continuous representation of \(G\) in \(L^p(T, \nu)\) for \(1 \leqslant p < +\infty\) . (One shows, as in the proof of Prop. 9 of §2, No.~5, that \(s \mapsto \gamma_{\chi}(s)\) is a representation of \(G\) by endomorphisms of \(L^p(T, \nu)\) . One then observes that if \(f \in L^p(T, \nu)\) and \(h \in \mathcal{K}(G)\) , Lemma 1 of §4, No.~1 shows that \(h(s)f(s^{-1}x)\chi(s^{-1}, x)\) is \(p\) -th power integrable for \(\beta \otimes \nu\) . The proof is concluded as in c).)

\begin{enumerate}
\def\labelenumi{\arabic{enumi})}
\setcounter{enumi}{13}
\tightlist
\item
  \begin{enumerate}
  \def\labelenumii{\alph{enumii})}
  \tightlist
  \item
    Let \(G\) be a locally compact group, \(f\) a lower semi-continuous positive function on \(G\) , \(\mu\) a positive measure on \(G\) . Show that the function
  \end{enumerate}
\end{enumerate}

\[
x \mapsto \int^ {*} f (s ^ {- 1} x) d \mu (s) = g (x)
\]

is lower semi-continuous on G (cf.~Ch. IV, §1, No.~1, Th. 1); for \(\mu\) and \(f\) to be convolvable, it is necessary and sufficient that \(g\) be integrable for a left Haar measure on G.

\begin{enumerate}
\def\labelenumi{\alph{enumi})}
\setcounter{enumi}{1}
\tightlist
\item
  On the group \(\mathbf{G} = \mathbf{R} \times \mathbf{R}\) , let \(\mu\) be the measure \(\varepsilon_0 \otimes \lambda\) , where \(\lambda\) is Lebesgue measure on \(\mathbf{R}\) ; let \(f(x, y) = (1 - |xy - 2|)^+\) ; show that the function
\end{enumerate}

\[
g (x, y) = \int f (x - s, y - t) d \mu (s, t)
\]

is everywhere finite on G, but is not continuous and is not integrable for Lebesgue measure on G.

¶ 15) Let G be a locally compact group, β a left Haar measure on G; by an abuse of language, in what follows one identifies β-integrable numerical functions with their classes in L¹(G,β); same abuse for the Lᵖ(G,β).

\begin{enumerate}
\def\labelenumi{\alph{enumi})}
\item
  Let \(A\) be a continuous endomorphism of the Banach space \(\mathrm{L}^1(\mathrm{G},\beta)\) such that, for every \(s \in \mathrm{G}\) , one has \(A(f*\varepsilon_s) = A(f)*\varepsilon_s\) for all \(f \in \mathrm{L}^1(\mathrm{G},\beta)\) . Show that for every function \(g \in \mathcal{K}(\mathrm{G})\) , one then has \(A(f*g) = A(f)*g\) (observe that \(s \mapsto g(s)A(f*\varepsilon_s)\) is a \(\beta\) -integrable mapping of \(\mathrm{G}\) into \(\mathrm{L}^1(\mathrm{G},\beta)\) ); converse. From this, deduce that there exists one and only one bounded measure \(\mu\) on \(\mathrm{G}\) such that \(A(f) = \mu *f\) for all \(f \in \mathrm{L}^1(\mathrm{G},\beta)\) , and that \(\| A\| = \| \mu \|\) . (With the notations of Prop. 19 of No.~7, consider the limit of \(A(f_{\mathrm{V}}*g)\) with respect to an ultrafilter finer than the section filter of \(\mathfrak{B}\) , making use of the compactness of the unit ball of \(\mathcal{M}^1(\mathrm{G})\) for the weak topology \(\sigma (\mathcal{M}^1(\mathrm{G}),\overline{\mathcal{K}(\mathrm{G})})\) .)
\item
  Let \(\mu\) be a bounded measure on G; show directly that the norm of the continuous endomorphism \(\pmb{\gamma}(\mu): f \mapsto \mu * f\) of \(\mathrm{L}^{\infty}(\mathrm{G},\beta)\) is equal to \(\| \mu \|\) . (Reduce to the case that \(\mu\) has compact support and has a continuous density with respect to \(|\mu|\) .) From this, deduce anew that the continuous endomorphism \(\pmb{\gamma}(\mu): f \mapsto \mu * f\) of \(\mathrm{L}^{1}(\mathrm{G},\beta)\) has norm equal to \(\| \mu \|\) .
\item
  Assume that G is compact and \(\mu\) is positive. Show that for \(1 < p < +\infty\) , the norm of the continuous endomorphism \(\gamma(\mu): f \mapsto \mu * f\) of \(\mathrm{L}^{p}(\mathrm{G}, \beta)\) is equal to \(\|\mu\|\) .
\item
  Take for G the cyclic group of order 3. Give an example of a measure \(\mu\) on G such that the norm of the endomorphism \(\pmb{\gamma}(\mu)\) of \(\mathrm{L}^p (\mathrm{G},\beta)\) is strictly less than \(\| \mu \|\) for \(1 < p < +\infty\) .
\end{enumerate}

¶ 16) Notations and conventions are those of Exer. 15.

\begin{enumerate}
\def\labelenumi{\alph{enumi})}
\tightlist
\item
  Show that, for a bounded measure \(\mu\) on \(\mathbf{G}\) to be such that \(\| \mu * f \|_1 = \| f \|_1\) for all \(f \in \mathrm{L}^1(\mathbf{G}, \beta)\) , it is necessary and sufficient that \(\mu\) be a point measure of norm 1. (Using the fact that the endomorphism \(\gamma(|\mu|)\) of \(\mathrm{L}^1(\mathbf{G}, \beta)\) has norm \(\| \mu \|\) (Exer. 15), show that one necessarily has, for every function \(f \in \mathcal{K}(\mathbf{G})\) ,
\end{enumerate}

\[
\left| \int f d \mu \right| = \int | f | d | \mu |;
\]

from this, first deduce that \(\mu = c|\mu|\) , where \(c\) is a constant of absolute value 1, then that \(\mu\) is a point measure.)

\begin{enumerate}
\def\labelenumi{\alph{enumi})}
\setcounter{enumi}{1}
\tightlist
\item
  Take for G the cyclic group of order 3. Give an example of a measure \(\mu\) on G, not a point measure, such that \(\| \mu * f \|_2 = \| f \|_2\) for every numerical function \(f\) defined on G.
\end{enumerate}

\begin{enumerate}
\def\labelenumi{\arabic{enumi})}
\setcounter{enumi}{16}
\tightlist
\item
  Notations and conventions are those of Exer. 15.
\end{enumerate}

\begin{enumerate}
\def\labelenumi{\alph{enumi})}
\item
  Let \(\mu\) be a bounded measure on G; for the endomorphism \(\pmb{\gamma}(\mu): f \mapsto \mu * f\) of \(\mathrm{L}^p(\mathrm{G}, \beta)\) to be surjective ( \(1 \leqslant p < +\infty\) ), it is necessary and sufficient that there exist a \(c_p > 0\) such that \(\| \pmb{\mu}^{\vee} * g \|_q \geqslant c_p \| g \|_q\) for all \(g \in \mathrm{L}^q(\mathrm{G}, \beta)\) (where \(\frac{1}{p} + \frac{1}{q} = 1\) ) (cf.~TVS, IV, §4, No.~2, remarks following Cor. 3 of Th. 1).
\item
  If \(\mu\) is a measure with base \(\beta\) , and G is not discrete, show that \(\gamma(\mu)\) is never surjective for \(1 \leqslant p \leqslant +\infty\) (for \(p < +\infty\) , use \(a\) ), reducing to the case that the density of \(\mu\) with respect to \(\beta\) belongs to \(\mathcal{K}(\mathrm{G})\) ; for \(p = +\infty\) , argue directly by observing that for \(f \in \mathrm{L}^1(\mathrm{G}, \beta)\) and \(g \in \mathrm{L}^\infty(\mathrm{G}, \beta)\) , \(f * g\) is uniformly continuous for the right uniform structure).
\item
  If G is abelian, \(1 \leqslant p \leqslant 2\) and the endomorphism \(\pmb{\gamma}(\mu)\) of \(\mathrm{L}^p(\mathrm{G},\beta)\) is surjective, show that it is bijective (using regularization, show that if \(\pmb{\gamma}(\mu)\) is not injective then the endomorphism \(\pmb{\gamma}(\check{\mu})\) of \(\mathrm{L}^q(\mathrm{G},\beta)\) is not injective: one makes use of Ch. IV, §6, No.~5, Cor. of Prop. 4).
\item
  Take \(\mathbf{G} = \mathbf{Z}\) and \(\mu = \varepsilon_1 - \varepsilon_0\) ; show that \(\pmb{\gamma}(\mu)\) is injective in the \(\mathrm{L}^p (\mathbf{Z},\beta)\) for which \(p\neq +\infty\) , but not in \(\mathrm{L}^{\infty}(\mathbf{Z},\beta)\) , and it is not surjective for any value of \(p\) .
\end{enumerate}

¶ 18) Notations and conventions are those of Exer. 15.

\begin{enumerate}
\def\labelenumi{\alph{enumi})}
\item
  Let \(\mu\) be a bounded measure on \(G\) , whose support contains at least two distinct points. Show that there exists a compact set \(K\) , and a number \(k\) such that \(0 < k < 1\) , for which \(\| \varphi_{sK} \cdot \mu \| \leqslant k \| \mu \|\) for all \(s \in G\) . (If \(t, t'\) are two distinct points in the support of \(\mu\) , take for \(K\) a neighborhood of \(e\) sufficiently small that \(sK\) cannot intersect both \(tK\) and \(t'K\) .)
\item
  Let \(\mu\) be a bounded measure \(\geqslant 0\) on \(G\) , whose support contains at least two distinct points. Show that there exists a function \(f \in L^{\infty}(G, \beta)\) , not equivalent to a function \(\geqslant 0\) , such that \(\mu * f\) is \(\geqslant 0\) locally almost everywhere for \(\beta\) . (Use \(a\) ), taking \(f\) to be equal to \(-1\) on \(K\) and to a suitable positive constant on \(\mathbb{C}K\) .) If, in addition, the support of \(\mu\) is compact, there exists a function \(f\) having the preceding properties and having compact support.
\item
  Let \(\mu_{1},\mu_{2}\) be two nonzero, bounded positive measures on G, permutable for convolution, such that the support K of \(\mu_{1}\) is compact and contains \(e\) . Set \(\mu = \mu_{1} + \mu_{2}\) . Let V be a symmetric compact neighborhood of \(e\) containing K, \(g\) the function equal to \(\mu * \varphi_{\mathrm{V}}\) on \(\mathbb{C}\mathrm{V}\) , and to 0 on V; show that the function \(\mu * (g - (\mu_1 * \varphi_{\mathrm{V}}))\) is \(\geqslant 0\) locally almost everywhere in \(\mathbb{C}(\mathrm{V}^2)\) . On the other hand, show that there exists a compact set H such that the function \(\mu_{2} * \varphi_{\mathrm{H}}\) is \(\geqslant 0\) almost everywhere in \(\mathrm{V}^2\) .
\item
  Deduce from c) that if \(\mu\) is a bounded positive measure on G whose support contains at least two distinct points, there exists a function f that belongs to every \(L^{p}(G,\beta)\) ( \(1 \leqslant p \leqslant +\infty\) ), is not equivalent to a function \(\geqslant 0\) , and is such that \(\mu * f\) is \(\geqslant 0\) locally almost everywhere, when one makes in addition one of the following hypotheses: \(\alpha\) ) G is abelian; \(\beta\) ) there is a point \(a \in G\) such that \(\mu(\{a\}) > 0\) . (Suitably decompose \(\mu\) as a sum of two permutable measures \(\geqslant 0\) .)
\end{enumerate}

¶ 19) Notations and conventions are those of Exer. 15.

\begin{enumerate}
\def\labelenumi{\alph{enumi})}
\item
  Show that if a positive bounded measure \(\mu\) on \(G\) is such that for every function \(f \in L^{p}(G, \beta)\) ( \(p\) given, \(1 \leqslant p < \infty\) ) the relation \(\mu * f \geqslant 0\) locally almost everywhere implies \(f \geqslant 0\) locally almost everywhere, then one can conclude that \(\mu\) is a point measure in each of the following cases: \(\alpha)\) \(G\) is compact; \(\beta)\) \(G\) is discrete; \(\gamma)\) \(G\) is abelian (use Exer. 18) (*). For \(p = +\infty\) , the same conclusion is valid without a supplementary hypothesis on \(G\) .
\item
  Let \(\mu\) be a positive bounded measure on \(G\) such that: \(1^{\circ} \pmb{\gamma}(\mu)\) is a surjective endomorphism of \(L^{1}(G, \beta)\) ; \(2^{\circ}\) for every function \(f \in L^{1}(G, \beta)\) , the relation \(\mu * f \geqslant 0\) locally almost everywhere implies \(f \geqslant 0\) locally almost everywhere. Show that \(\mu\) is then a point measure. (Observe that for every function \(g \in L^{\infty}(G, \beta)\) , the relation \(\stackrel{\vee}{\mu} * g \geqslant 0\) locally almost everywhere implies \(g \geqslant 0\) locally almost everywhere.)
\end{enumerate}

\begin{enumerate}
\def\labelenumi{\arabic{enumi})}
\setcounter{enumi}{19}
\tightlist
\item
  Let \(G\) be a locally compact group, \(\beta\) a left Haar measure on \(G\) .
\end{enumerate}

\begin{enumerate}
\def\labelenumi{\alph{enumi})}
\item
  Let \(f\) be a bounded numerical function on \(G\) , uniformly continuous for the left uniform structure. Show that for every bounded measure \(\mu\) on \(G\) , \(\mu * f\) is uniformly continuous for the left uniform structure; moreover, with the notations of No.~7, Prop. 19, \(f\) is the limit, for the topology of uniform convergence in \(G\) , of the functions \(f_V * f\) with respect to the section filter of \(\mathfrak{B}\) .
\item
  Take \(\mathbf{G} = \mathbf{T}\) ; give an example of a continuous function \(h\) on \(\mathbf{T}\) that is not of the form \(f*g\) , where \(f\) and \(g\) belong to \(\mathrm{L}^2 (\mathbf{T},\beta)\) . (Use Exers. 15 b) and 16 of Ch. IV, §6.)
\end{enumerate}

¶ 21) Let G be a locally compact group, β a left Haar measure on G; one canonically identifies L¹(G,β) with a subspace of M¹(G).

\begin{enumerate}
\def\labelenumi{\alph{enumi})}
\item
  With the notations of §3, Exer. 11, let A be a relatively compact subset of \(\mathcal{M}_{\mathrm{III}}\) , B a subset of \(\mathrm{L}^1 (\mathrm{G},\beta)\) , relatively compact in \(\mathcal{M}_{\mathrm{IV}}\) . Show that the restriction to \(\mathrm{A}\times \mathrm{B}\) of the mapping \((\mu ,\nu)\mapsto \mu *\nu\) of \(\mathcal{M}_{\mathrm{III}}\times \mathcal{M}_{\mathrm{IV}}\) into \(\mathcal{M}_{\mathrm{IV}}\) is continuous. (First prove that if A and B are compact, the image of \(\mathrm{A}\times \mathrm{B}\) under this mapping is compact; use Exer. 10 of §3, as well as the criterion \(\alpha\) ) of Ch. V, §5, Exer. 15. To show next that \((\mu ,\nu)\mapsto \mu *\nu\) is continuous, make use of Exer. 11 c) of §3.)
\item
  Let \(\mathcal{U}_s^\infty (\mathrm{G})\) be the set of bounded numerical functions on \(\mathbf{G}\) uniformly continuous for the left uniform structure. Denote by \(\mathcal{T}_{\mathrm{II}}\) the topology \(\sigma (\mathcal{M}^1 (\mathrm{G}),\mathcal{U}_s^\infty (\mathrm{G}))\) and by \(\mathcal{M}_{\mathrm{II}}\) the space \(\mathcal{M}^1 (\mathrm{G})\) equipped with \(\mathcal{T}_{\mathrm{II}}\) . Taking \(\mathbf{G} = \mathbf{R}\) , give an example of a sequence \((\mu_n)\) of measures on \(\mathbf{G}\) tending to 0 for \(\mathcal{T}_{\mathrm{II}}\) and a sequence \((f_n)\) of functions in \(\mathrm{L}^1 (\mathrm{G},\beta)\) , tending to 0 for \(\mathcal{T}_{\mathrm{IV}}\) (or, what comes to the same, for the topology \(\sigma (\mathrm{L}^1 (\mathrm{G},\beta),\mathrm{L}^\infty (\mathrm{G},\beta))\) ), such that the sequence \((\mu_n*f_n)\) does not tend to 0 for \(\mathcal{T}_{\mathrm{IV}}\) (take \(f_{n}(t) = \sin nt\) in the interval \([0,\pi ]\) , \(f_{n}(t) = 0\) elsewhere).
\item
  Let A be a relatively compact subset of \(\mathcal{M}_{\mathrm{II}}\) , B a subset of \(\mathrm{L}^1 (\mathrm{G},\beta)\) , relatively compact for the topology of the norm \(\| f\| _1\) . Show that the restriction to \(\mathbf{A}\times \mathbf{B}\) of the mapping \((\mu ,f)\mapsto \mu *\textit{f}\) of \(\mathcal{M}_{\mathrm{III}}\times \mathrm{L}^1 (\mathrm{G},\beta)\) into \(\mathrm{L}^1 (\mathrm{G},\beta)\) is continuous (L \(^1 (\mathrm{G},\beta)\) being equipped with its normed space topology). (Reduce to proving that for \(f\in \mathcal{K}(\mathrm{G})\) , the mapping \(\mu \mapsto \mu *\textit{f}\) of \(\mathcal{M}_{\mathrm{II}}\) into \(\mathrm{L}^1 (\mathrm{G},\beta)\) is continuous in A. Using the fact that for \(g\in \mathrm{L}^{\infty}(\mathrm{G},\beta)\) and \(f\in \mathcal{K}(\mathrm{G})\) , \(g*\textit{f}\) is uniformly continuous for the left uniform structure, first show that the mapping \(\mu \mapsto \mu *\textit{f}\) of \(\mathcal{M}_{\mathrm{II}}\) into \(\mathcal{M}_{\mathrm{IV}}\) is continuous. Next, using Exer. 20 a), reduce to proving that if C is a subset of \(\mathrm{L}^1 (\mathrm{G},\beta)\) compact for the topology \(\sigma (\mathrm{L}^{1}(\mathrm{G},\beta),\mathrm{L}^{\infty}(\mathrm{G},\beta))\) and \(h\in \mathcal{K}(\mathbf{G})\) , the mapping \(\mu \mapsto \mu *h\) of C into \(\mathrm{L}^1 (\mathrm{G},\beta)\) is continuous when \(\mathrm{L}^1 (\mathrm{G},\beta)\) is equipped with its normed space topology. By the same reasoning as in \(a\) ), show that for this it suffices to prove that the image of C under this mapping is compact for the normed space topology. To this end, use Šmulian's theorem, the fact that C is cramped (§3, Exer. 10) and Lebesgue's theorem.)
\item
  Let A be a relatively compact subset of \(\mathcal{M}_{\mathrm{II}}\) containing 0, B a relatively compact subset of \(\mathcal{M}_{\mathrm{I}}\) ; show that for every \(\nu_0 \in \mathrm{B}\) , the restriction to \(\mathrm{A} \times \mathrm{B}\) of the mapping \((\mu, \nu) \mapsto \mu * \nu\) of \(\mathcal{M}_{\mathrm{II}} \times \mathcal{M}_{\mathrm{I}}\) into \(\mathcal{M}_{\mathrm{II}}\) is continuous at the point \((0, \nu_0)\) . (Make use of Exer. 20 a), and Exer. 21 c).
\item
  Let A be a bounded subset of \(\mathcal{M}^1 (\mathrm{G})\) , \(f\) a function in \(\mathrm{L}^1 (\mathrm{G},\beta)\) ; show that the restriction to A of the mapping \(\mu \mapsto \mu *f\) of \(\mathcal{M}_{\mathrm{III}}\) into \(\mathcal{M}_{\mathrm{IV}}\) is continuous.
\item
  Let \(\mathcal{M}_{+}^{1}(G)\) be the set of bounded positive measures on G. Show that if a subset A of \(\mathcal{M}_{+}^{1}(G)\) is compact for \(T_{II}\) , it is also compact for \(T_{III}\) . (Reduce to proving that A is a cramped set. Argue by contradiction, using the fact that for \(f \in \mathcal{K}(G)\) , the image of A under the mapping \(\mu \mapsto \mu * f\) is a cramped set by virtue of c).)
\end{enumerate}

\(g)\) Show that for \(\mathbf{G} = \mathbf{R}\) , the topologies induced on \(\mathcal{M}_{+}^{1}(\mathbf{G})\) by \(\mathcal{T}_{\mathrm{II}}\) and \(\mathcal{T}_{\mathrm{III}}\) are distinct.

¶ 22) Notations are those of Exer. 21 of §4 and Exer. 11 of §3.

\begin{enumerate}
\def\labelenumi{\alph{enumi})}
\item
  Let \((\mu_n)\) be a sequence of bounded measures on \(G\) . Assume that for every function \(f \in L^1(G, \beta)\) , the sequence \((\mu_n * f)\) converges to 0 in \(\mathcal{M}_I\) . Show that the sequence of norms \((\|\mu_n\|)\) is bounded (use the Banach-Steinhaus theorem for the family of mappings \(f \mapsto \mu_n * f\) of \(L^1(G, \beta)\) into itself, as well as Exer. 15). Deduce from this that the sequence \((\mu_n)\) tends to 0 in \(\mathcal{M}_I\) (use Exer. 20).
\item
  Let \((\mu_n)\) be a sequence of bounded measures on \(G\) such that the sequence \((\mu_n * f)\) tends vaguely to 0 for every function \(f \in L^1(G, \beta)\) ; show that the sequence \((\mu_n)\) tends vaguely to 0 (for every compact subset \(K\) of \(G\) , show, arguing as in \(a\) ), that the sequence \((|\mu_n|(K))\) is bounded). Give an example (with \(G = Z\) ) where the sequence \((\|\mu_n\|)\) is not bounded.
\item
  Let \((\mu_n)\) be a sequence of bounded measures on \(G\) such that for every function \(f \in L^1(G, \beta)\) , the sequence \((\mu_n * f)\) converges to 0 in \(\mathcal{M}_{II}\) ; show that the sequence \((\mu_n)\) converges to 0 in \(\mathcal{M}_{II}\) (make use of \(a\) )).
\end{enumerate}

\begin{enumerate}
\def\labelenumi{\arabic{enumi})}
\setcounter{enumi}{22}
\item
  Let \(G\) be a locally compact group, \(\beta\) a left Haar measure on \(G\) , \(f\) and \(g\) two positive locally \(\beta\) -integrable functions. Assume that \(f\) and \(g\) are convolvable and that one of these functions is zero on the complement of a countable union of compact sets. Show that \(f * g\) is equal locally almost everywhere to a lower semi-continuous function (use Prop. 15 of No.~5).
\item
  Show that the inequalities of Props. 12 and 15 of No.~5 cannot be improved by inserting a constant factor c \textless{} 1 in their right members.
\item
  Let G be a locally compact group, \(\beta\) a left Haar measure on G, \(\mu\) a positive measure on G. If A is a \(\mu\) -integrable set and B is a Borel set in G, the function
\end{enumerate}

\[
u: s \mapsto \mu (\mathrm{A} \cap s \mathrm{B})
\]

is \(\beta\) -measurable in G; if, in addition, \(B^{-1}\) is \(\beta\) -integrable, then so is \(u\) , and

\[
\int \mu (\mathrm{A} \cap s \mathrm{B}) d \beta (s) = \mu (\mathrm{A}) \beta (\mathrm{B} ^ {- 1}).
\]

Give an example of a measure \(\mu\) such that \(s \mapsto \mu(A \cap sB)\) is not continuous; if \(\mu\) is a measure with base \(\beta\) , and if A is \(\mu\) -integrable and B is a Borel set, then the function \(s \mapsto \mu(A \cap sB)\) is continuous on G.

¶ 26) Let G be a locally compact group, β a left Haar measure on G. For a subset H of L\^{}p(G,β) (1 ≤ p \textless{} +∞) to be relatively compact (for the topology of convergence in mean of order p), it is necessary and sufficient that the following conditions be satisfied: 1° H is bounded in L\^{}p(G,β); 2° for every ε \textgreater{} 0, there exists a compact subset K of G such that \textbar\textbar fφ\_G-K\textbar\textbar\_p ≤ ε for every function f ∈ H; 3° for every ε \textgreater{} 0, there exists a neighborhood V of e in G such that \textbar\textbar γ(s)f - f\textbar\textbar\_p ≤ ε for all f ∈ H and s ∈ V. (To prove that the conditions are sufficient, observe that if g ∈ K(G) and if L is a compact subset of G, then the image, under the mapping f ↦ g * f, of the set φ\_L · H, is an equicontinuous subset of K(G).)

¶ 27) Let G be a locally compact group, β a left Haar measure on G. If G is not reduced to e, then L¹(G,β) is an algebra (for convolution) admitting divisors of zero other than 0. One may proceed as follows to form two nonzero elements f, g of L¹(G,β) such that f * g = 0:

\(1^{\circ}\) The case that \(\mathbf{G}\) contains a compact subgroup \(\mathbf{H}\) , not reduced to \(e\) , in which \(\Delta(x)=1\) . Take \(f\) to be the characteristic function of a set \(\mathbf{A}\) , and \(g\) to be a difference \(\varphi_{s\mathbf{B}}-\varphi_{\mathbf{B}}\) of two characteristic functions, with \(\mathbf{A}\) and \(\mathbf{B}\) suitably chosen.

\(2^{\circ}\) The case that \(\mathbf{G} = \mathbf{Z}\) . Show that one can then take

\[
f (n) = \frac {1}{2 n - 1} - \frac {1}{2 n + 1}
\]

for all \(n \in \mathbf{Z}\) , and \(g(n) = f(-n)\) .

\(3^{\circ}\) The general case. First prove that there exists an \(a \neq e\) in G such that \(\Delta(a) = 1\) . The closure H in G of the subgroup generated by \(a\) is then either a compact subgroup or a subgroup isomorphic to Z (GT, V, §1, Exer. 2). In the first case, use the result of \(1^{\circ}\) ; in the second, take

\[
f (t) = \sum_ {n = - \infty} ^ {+ \infty} \alpha_ {n} \varphi_ {\mathrm{U} a ^ {- n}} (t), \quad g (t) = \sum_ {n = - \infty} ^ {+ \infty} \beta_ {n} \varphi_ {a ^ {n} \mathrm{U}} (t),
\]

suitably choosing U and the sequences \((\alpha_{n})\) , \((\beta_{n})\) with the help of \(2^{\circ}\) .

\begin{enumerate}
\def\labelenumi{\arabic{enumi})}
\setcounter{enumi}{27}
\tightlist
\item
  Let \(G_1, G_2\) be two locally compact groups, \(\beta_1, \beta_2\) left Haar measures on \(G_1, G_2\) , and \(A_1\) (resp. \(A_2\) ) the topological algebra (over \(R\) ) \(L^1(G_1, \beta_1)\) (resp. \(L^1(G_2, \beta_2)\) ). Let \(T\) be an algebra isomorphism of \(A_1\) onto \(A_2\) , such that the relation \(f \geqslant 0\) almost everywhere is equivalent to \(T(f) \geqslant 0\) almost everywhere.
\end{enumerate}

\begin{enumerate}
\def\labelenumi{\alph{enumi})}
\item
  Show that \(\mathbf{T}\) is an isomorphism of topological algebras. (First note that if a decreasing sequence \((f_n)\) of elements of \(\mathbf{A}_1\) tends to 0 almost everywhere, then the same is true of the sequence of the \(\mathrm{T}(f_n)\) ; deduce from this that for every sequence \((f_n)\) that is bounded above in \(\mathbf{A}_1\) , \(\mathrm{T}(\sup (f_n)) = \sup (\mathrm{T}(f_n))\) , and conclude with the help of Fatou's lemma.)
\item
  Show that T may be extended, in only one way, to an isomorphism of the topological algebra \(\mathcal{M}^{1}(G_{1})\) onto the topological algebra \(\mathcal{M}^{1}(G_{2})\) , and that there exist a topological isomorphism u of \(G_{1}\) onto \(G_{2}\) and a continuous homomorphism \(\chi\) of \(G_{1}\) into \(R_{+}^{*}\) such that \(\mathrm{T}(\varepsilon_{s}) = \chi(s)\varepsilon_{u(s)}\) (make use of Exer. 15 a) and Exer. 19 b); to prove that u and \(\chi\) are continuous, observe that T defines an isomorphism of the algebra \(\mathcal{L}(A_{1})\) of continuous endomorphisms of the topological vector space \(A_{1}\) onto the analogous algebra \(\mathcal{L}(A_{2})\) , and that this isomorphism is bicontinuous for the topology of pointwise convergence; on the other hand, observe that \(s \mapsto \delta(s^{-1})\) is an isomorphism of \(G_{1}\) onto a multiplicative subgroup of \(\mathcal{L}(A_{1})\) .
\item
  Deduce from \(b)\) that, for every \(f \in \mathbf{A}_1\) ,
\end{enumerate}

\[
(T (f)) (t) = \chi (u ^ {- 1} (t)) f (u ^ {- 1} (t))
\]

for all \(t \in \mathbf{G}_2\) .

Recall (A, III, §2, Exer. 9) that there exist finite groups \(\mathbf{G}_1, \mathbf{G}_2\) such that the algebras \(\mathrm{L}^1(\mathrm{G}_1, \beta_1)\) and \(\mathrm{L}^1(\mathrm{G}_2, \beta_2)\) are isomorphic without \(\mathbf{G}_1\) and \(\mathbf{G}_2\) being so.

\exercisesection{}

\begin{enumerate}
\def\labelenumi{\arabic{enumi})}
\tightlist
\item
  Let X be a locally compact space, \(\mathfrak{F}\) (or \(\mathfrak{F}(X)\) ) the set of closed subsets of X. For every compact subset K of X, every compact neighborhood L of K, and every entourage V of the unique uniform structure of L, let \(Q(K,L,V)\) be the set of pairs (A,B) of elements of \(\mathfrak{F}\) satisfying the two conditions
\end{enumerate}

\[
\mathrm{A} \cap \mathrm{K} \subset \mathrm{V} (\mathrm{B} \cap \mathrm{L}) \quad \text { and } \quad \mathrm{B} \cap \mathrm{K} \subset \mathrm{V} (\mathrm{A} \cap \mathrm{L}).
\]

\begin{enumerate}
\def\labelenumi{\alph{enumi})}
\item
  Show that the sets \(\mathbf{Q}(\mathbf{K},\mathbf{L},\mathbf{V})\) form a fundamental system of entourages for a Hausdorff uniform structure on \(\mathfrak{F}\) .
\item
  Let \(\mathcal{U}\) be a uniform structure compatible with the topology of X. For every compact subset K of X and every entourage W of \(\mathcal{U}\) , let P(K, W) be the set of pairs (A, B) of \(\mathfrak{F}\) satisfying the two conditions
\end{enumerate}

\[
\mathrm{A} \cap \mathrm{K} \subset \mathrm{W} (\mathrm{B}) \quad \text { and } \quad \mathrm{B} \cap \mathrm{K} \subset \mathrm{W} (\mathrm{A}).
\]

Show that the sets \(\mathbf{P}(\mathbf{K},\mathbf{W})\) form a fundamental system of entourages for the uniform structure defined in \(a\) ).

\begin{enumerate}
\def\labelenumi{\alph{enumi})}
\setcounter{enumi}{2}
\tightlist
\item
  Let \(X'\) be the Alexandroff compactification of X, \(\omega\) the point at infinity of \(X'\) ; for every closed subset A of X, let \(A' = A \cup \{\omega\}\) . Show that the mapping \(A \mapsto A'\) is an isomorphism of the uniform space \(\mathfrak{F}(X)\) defined in a) onto the subspace of \(\mathfrak{F}(X')\) formed by the closed subsets containing \(\omega\) (make use of b)). Deduce from this that \(\mathfrak{F}(X)\) is compact (cf.~GT, II, §4, Exer. 15).
\end{enumerate}

\begin{enumerate}
\def\labelenumi{\arabic{enumi})}
\setcounter{enumi}{1}
\item
  Let \(G\) be a locally compact group, \(\mathfrak{F}(G)\) the uniform space of closed subsets of \(G\) (defined in No.~6 or in Exer. 1). Prove directly that the set \(\Sigma\) of closed subgroups of \(G\) is closed in \(\mathfrak{F}(G)\) (without using Prop. 3 of No.~1).
\item
  Let \(G\) be a locally compact group, \(\chi\) a continuous representation of \(G\) in \(\mathbf{R}_+^*\) . Generalize Props. 4, 5 and 6 to the set \(\Gamma_{\chi}\) of measures \(\alpha \in \Gamma\) such that the restriction of \(\chi\) to \(H_{\alpha}\) is the modulus of \(\alpha\) . (Replace the measure \(\mu\) by \(\chi \cdot \mu\) and consider the quotient measures \((\chi \cdot \mu) / \overset{\vee}{\alpha}\) .)
\item
  Let G be a locally compact group that is not generated by any compact subset of G, and let H be a discrete subgroup of G such that G/H is compact. Show that H cannot be generated by a finite number of elements, but the normalized Haar measure \(\alpha_{0}\) of H belongs to the closure in \(\Gamma^{0}\) of the set of \(\alpha\) such that \(\|\mu_{\alpha}\| = +\infty\) (consider the normalized Haar measures of the subgroups of H that are generated by a finite number of elements).
\item
  Let G be the discrete group that is the direct sum of an infinite sequence \((\mathrm{G}_{n})\) of subgroups with two elements, and let \(p_{n}\) be the projection of G onto \(G_{n}\) . Set \(\mathrm{H}_{n} = \overline{p_{n}}(e)\) , and let \(\alpha_{n}\) be the normalized Haar measure of \(H_{n}\) . If \(\alpha\) is the normalized Haar measure of G, show that in \(\Gamma_{c}^{0}\) the sequence \((\alpha_{n})\) converges to \(\alpha\) but that \(\|\mu_{\alpha_{n}}\| = 2\) and \(\|\mu_{\alpha}\| = 1\) .
\item
  Let G be a locally compact group that does not satisfy the condition (L) of No.~4.
\end{enumerate}

\begin{enumerate}
\def\labelenumi{\alph{enumi})}
\item
  Show (with the notations of No.~4) that the canonical bijection of N onto D is not a homeomorphism. (For every neighborhood W of e in G, let A(W) be the set of finite subgroups \(\neq\{e\}\) contained in W; show that the A(W) form the base of a filter that converges in D to the subgroup \(\{e\}\) , but that if \(\alpha_{H}\) denotes the normalized Haar measure of H, the measures \(\alpha_{H}\) do not converge to \(\varepsilon_{e}\) .)
\item
  If in addition G is metrizable, show that there exists a compact subset B of the space D of discrete subgroups of G, that is not contained in any of the sets \(D_{U}\) defined in Th. 1 of No.~3.
\end{enumerate}

\specialchapter{Historical Note (Chapters VII and VIII)}
